\chapter{Introduction}

This chapter introduces the research problem, motivation, objectives, and scope of the thesis. It first outlines the background and practical significance of ransomware detection on Linux systems, then defines the key problems addressed in this work. The chapter next presents the research objectives, research questions, and principal contributions, and concludes with an overview of the organization of the thesis.

\section{Background and Motivation}

Ransomware has emerged as one of the most destructive categories of malware confronting organizations worldwide. Unlike data-theft malware that exfiltrates information covertly, ransomware encrypts victim files and demands payment for decryption keys, inflicting immediate and quantifiable harm. According to IBM Security's Cost of a Data Breach Report 2025, the average cost of a ransomware breach reached USD 5.08 million per incident, while the global average across all breach types stood at USD 4.44 million (with the United States reaching an all-time high of USD 10.22 million) \cite{ref1}. High-profile incidents targeting critical infrastructure, healthcare systems, and government entities have demonstrated that ransomware poses not only severe financial risks but also significant public-safety consequences.

Linux-based systems have become increasingly prominent ransomware targets \cite{ref2}. The widespread adoption of Linux in cloud infrastructure, enterprise servers, and virtualization platforms (particularly VMware ESXi hypervisors \cite{ref3}) has made it an attractive attack surface. Modern ransomware families such as Akira \cite{ref4}, BlackBasta \cite{ref5}, IceFire \cite{ref6}, and LockBit \cite{ref7} have developed Linux-specific variants engineered to target Linux filesystem semantics. Because Linux servers frequently host high-value data repositories, a single successful ransomware execution can encrypt terabytes of data within minutes, with potentially catastrophic consequences for business continuity.

Existing ransomware detection approaches on Linux fall into two broad categories: signature-based detection and behavior-based detection. Signature-based methods rely on cryptographic hashes or byte-pattern matching and are rendered ineffective by minor code obfuscation or recompilation. Behavior-based methods analyze process activity patterns but many rely on user-space monitoring tools that introduce significant performance overhead, or require kernel module modifications that compromise compatibility with modern Linux kernels. Furthermore, most published detection systems were designed and evaluated against Windows ransomware behavior and do not generalize to Linux-specific encryption patterns. Detection approaches that rely on raw continuous features (such as absolute file sizes and wall-clock timestamps) are additionally susceptible to overfitting to the magnitude characteristics of known families and generalize poorly to novel strains whose I/\allowbreak{}O volume profiles differ significantly from training distributions. A candidate representation should therefore preserve relative behavioral patterns while reducing dependence on absolute magnitude.

Ransomware families increasingly employ evasive encryption strategies, such as encrypting only partial regions of target files, that deliberately mimic the I/\allowbreak{}O patterns of legitimate workloads, causing volume-based and entropy-based detectors to fail \cite{ref15}. These behaviors motivate evaluating temporal I/\allowbreak{}O sequences alongside aggregate statistics rather than assuming that either representation is universally sufficient.

The extended Berkeley Packet Filter (eBPF) subsystem offers a compelling foundation for Linux security monitoring. eBPF programs execute in a sandboxed kernel environment with minimal processing overhead, providing access to file-system events at the Virtual File System (VFS) layer without requiring kernel modification or loadable module insertion \cite{ref10}. This enables real-time, per-process event capture with sub-microsecond timestamping precision (via the kernel's nanosecond-resolution clock), offering lower overhead than user-space audit frameworks such as auditd or inotify \cite{ref10}, \cite{ref11}, \cite{ref40}. Despite these advantages, systematic collection and analysis of VFS-level I/\allowbreak{}O behavioral traces across multiple ransomware families, and their use as training data for stateful sequence models, remains underexplored in the research literature.

Certain in-place encryption implementations produce Write-after-Read (WaR) transitions, in which a process reads a byte range and subsequently writes to an overlapping range \cite{ref9}. The visibility of this pattern depends on the encryption implementation and the instrumented I/\allowbreak{}O path. WaR transitions also occur in legitimate workloads such as database compaction, log rotation, and software compilation; therefore, WaR is evaluated as one contextual sequence feature rather than as a universal indicator of ransomware.

\section{Problem Statement}

Despite significant research interest in ransomware detection, three fundamental problems remain inadequately addressed in the context of Linux system security:

1) Evasive Encryption Strategies:  Modern ransomware families employ selective and partial-file encryption strategies specifically designed to evade I/\allowbreak{}O-volume-based detectors \cite{ref15}. Certain families encrypt only fixed-size partial regions of each file and deliberately avoid overwriting file headers, producing event streams that closely resemble legitimate database compaction or log-rotation workloads. These evasive patterns cause threshold-based and entropy-based detectors to fail \cite{ref13}, \cite{ref14}; moreover, adversarial manipulation of behavioral parameters such as thread count and encryption ratio can further degrade AI-based detection \cite{ref8}. Such evasion techniques simultaneously remain invisible to detectors that rely solely on aggregate write volume or file-count statistics.

2) Lack of Stateful, Streaming Detection:  Ransomware behavior can evolve over time: an encryption campaign may begin slowly, accelerate, and then pause. Existing approaches commonly aggregate events into fixed-duration windows or extract static feature vectors from bounded event buffers \cite{ref16}. Resetting state between windows can alter sensitivity to behavior distributed across time, but the effect is behavior dependent. The ablation in Appendix F.4 found lower aggregate false positives and higher mean F1 with retained state, whereas resetting state improved Akira MDR. Cross-chunk state retention is therefore treated as an evaluated streaming design choice rather than as universally superior.

3) Limited Generalization to Unseen Families: Ransomware families continuously evolve. A detection model trained only on known families provides weak assurance against novel strains. Existing learning-based approaches rarely evaluate cross-family generalization rigorously for Linux ransomware \cite{ref17}, \cite{ref18} (for example, via a leave-one-family-out protocol that trains on all families except one and tests on the withheld family, repeated for each family), and those that do typically report aggregate metrics that obscure per-family weaknesses. There is currently no standardized evaluation benchmark for Linux ransomware that measures both detection latency and data-loss severity at the process level.

\section{Research Objectives}

The main objective of This thematic paper is to determine whether file I/\allowbreak{}O behavior captured at the Linux VFS layer using eBPF provides a sufficient signal for real-time ransomware detection that generalizes to previously unseen ransomware families. In support of this objective, the thesis pursues three sub-objectives:

(O1)	Design an eBPF-based detection system that captures file I/\allowbreak{}O behavior at the Linux VFS layer and classifies processes in real time.

(O2)	Evaluate whether a detector trained on known ransomware families generalizes to families withheld entirely from training.

(O3)	Investigate a two-stage refinement pipeline for suppressing false-positive alerts under realistic mixed workloads, and characterize its trade-offs against detection retention.

\section{Research Questions}

The thesis is guided by the following research questions:

(RQ1)	Can VFS-level file I/\allowbreak{}O operations serve as a sufficient behavioral signal for real-time ransomware detection on Linux?

(RQ2)	Does a ransomware detector trained on known families generalize to unseen families that employ different encryption strategies?

(RQ3)	How much data must a ransomware process corrupt before a VFS-based detector can raise an alert, and how does this vary across families with different behavioral profiles?

(RQ4)	Does a compact discrete tokenization of VFS I/\allowbreak{}O event streams capture sufficient behavioral signal to distinguish ransomware from benign processes without continuous feature normalization?

(RQ5)	To what extent can a two-stage detection pipeline (combining a streaming sequence classifier with a rule-assisted behavioral scoring stage) suppress false-positive alerts in a real-time deployment scenario, and what are its trade-offs with respect to missed detections?

\section{Scope and Research Contributions}

This research focuses on file-system ransomware targeting Linux server environments, using eBPF-based monitoring at the kernel VFS layer to capture real-time file I/\allowbreak{}O behavior across all active processes. The study encompasses fourteen Linux ransomware families, a purpose-built multi-family VFS I/\allowbreak{}O dataset, a compact event tokenization scheme, a stateful recurrent detection model, and a two-stage behavioral scoring pipeline, all evaluated under a rigorous cross-family generalization protocol. End-to-end live deployment is demonstrated as a proof-of-concept rather than as the primary evaluation; network-propagating ransomware variants (worms) and fileless ransomware that produce no file I/\allowbreak{}O are outside the scope of this study.

This thematic paper treats Linux ransomware detection primarily as a sensing problem rather than a modeling one. The guiding question is whether file-system behavior observable at a deployable kernel layer carries enough signal to flag an encrypting process early --- before meaningful data loss --- and at what coverage cost on a modern Linux I/\allowbreak{}O stack. The four contributions follow from that framing.

C1 --- A VFS-layer eBPF monitoring methodology and a multi-family Linux dataset. An eBPF sensor captures per-process read, write, and rename behavior at the Virtual File System layer using kprobes without requiring kernel modification; the evaluated deployment used the eBPF-related privileges described in Section 3.1.1. The sensor supported the systematic collection and quality audit of 682 Parquet recordings: 494 malware-session recordings across fourteen Linux ransomware families and 188 standalone benign-workload recordings. Direct counting of rows with Label = 1 in the malware-session recordings yielded approximately 598.6 million malware-labeled VFS I/\allowbreak{}O events. This infrastructure is the foundation on which the remaining contributions rest, and it directly addresses the absence of Linux-generated data noted by the CLEAR framework.

C2 --- Evidence of Discriminative Signal in the VFS Representation. The aggregate-feature baselines and the stateful \textmu{}LSTM achieved similar mean AUC values in the reported evaluation. However, differences in MDR indicate that the classifiers did not provide equivalent detection performance across all operational metrics. Furthermore, the comparison did not fully control the decision point, calibration procedure, stopping rule, and online memory requirements across the evaluated models. The results therefore suggest that the VFS representation carries strong discriminative signal and that classifier choice was not the dominant determinant of mean AUC in this experiment. They do not establish that the representation alone determines detection quality. The controlled SougoLock tokenization ablation provides additional evidence that insufficient representation resolution can cause detection failure, but this finding is limited to the evaluated fold and configuration.

C3 --- A characterization of what VFS-layer detection can and cannot see. A data-driven analysis of the fourteen families' I/\allowbreak{}O fingerprints --- write-size distributions, RENAME rates, opcode-transition patterns, and throughput profiles --- separates them into encryption-dominant, rename-dominant, and bimodal behavioral groups. The same analysis maps the structural limits of the approach: log-spam dilution, tokenizer offset-resolution limits, and most consequentially the I/\allowbreak{}O paths (io\_\allowbreak{}uring, pwritev, O\_\allowbreak{}DIRECT) that bypass the VFS hooks and are already exercised by families within this dataset. This coverage boundary is treated as a finding about the monitoring layer, not merely a limitation of the prototype (Sections 4.5, 5.3).

C4 --- A compact streaming detector and an evaluation calibrated to its evidence. As a deployability design point rather than an accuracy claim, the 182,849-parameter stateful \textmu{}LSTM consumes two 10-bit token IDs per accepted VFS event and maintains a fixed-size recurrent hidden state per process. Isolated model inference achieved sub-millisecond latency on the tested Windows CPU; this result supports a compact streaming design point but does not establish end-to-end continuous on-host performance. The detector is evaluated under a leave-one-family-out protocol with domain-specific metrics for missed detections and data exposure before the first alert. Results are reported separately for the five multi-binary folds, which test cross-family generalization in the intended sense, and the nine single-binary folds, for which the protocol reduces to a leave-one-binary-out experiment (Section 4.3), so the generalization claim is matched to what the dataset supports.

Taken together, the baseline comparison and the controlled tokenization ablation suggest that the Linux VFS representation carries substantial discriminative signal. However, these experiments do not establish that representation, rather than classifier choice, is the sole determinant of detection quality. The compared models differed in MDR, and equivalence was not established for calibration, early-detection behavior, memory requirements, or online deployability.

To make the structure of the research explicit, Table 1.1 maps each research objective to the research questions that investigate it and to the resulting contributions.

\begin{table}[htbp]
\centering
\small
\caption{Mapping of research objectives, research questions, and contributions.}\label{tab:1-1}
\begin{tabularx}{\textwidth}{XlX}
\hline
\textbf{Objective} & \textbf{Research Questions} & \textbf{Contributions} \\
\hline
O1: Design an eBPF-based VFS monitoring system that classifies processes in real time. & RQ1, RQ4 & C1, C4 \\
O2: Evaluate generalization to ransomware families withheld entirely from training. & RQ2, RQ3 & C2, C3 \\
O3: Investigate false-positive suppression and its trade-offs through a two-stage pipeline. & RQ5 & Empirical FP-reduction evidence (Chapter 4) \\
\hline
\end{tabularx}
\end{table}

\section{Thesis Organization}

The remainder of This thematic paper is organized into four further chapters that progressively develop, evaluate, and reflect on the proposed Linux ransomware detection system.

Chapter 2 reviews the background and related literature. It presents the theoretical foundations of ransomware behavior and taxonomy, surveys eBPF kernel instrumentation and the Linux VFS subsystem, summarizes LSTM-based sequence modeling as it applies to system-call and I/\allowbreak{}O event streams, and positions the present work relative to existing behavior-based ransomware detection approaches.

Chapter 3 describes the methodology. It details the eBPF VFS sensor design, the dataset collection and quality-audit procedures, the two-token event encoding using 10-bit token IDs, the stateful \textmu{}LSTM architecture, the training regime based on Truncated Backpropagation Through Time (TBPTT), the two-stage behavioral scoring pipeline, and the evaluation metrics including the domain-specific Missed Detection Rate (MDR) and MBD3, defined as the 75th percentile of cumulative WRITE bytes before first detection among detected malware processes.

Chapter 4 presents the experimental results. It opens with the evaluation setup and dataset partitions, reports leave-one-family-out (LOFO) detection performance across the fourteen Linux ransomware families with per-family MDR, MBD3, AUC, and confusion-matrix figures, characterizes the families through three behavioral groups, and reports isolated model-inference latency, throughput, and memory measurements from a separately trained prototype deployment. It then presents a three-part error analysis of cross-cutting failure modes (log-spam dilution, sensor blind-spots, and token-representation limits), the two-stage pipeline's false-positive suppression, and a live mixed-workload validation, concluding with a discussion that interprets the results against the contributions stated in Chapter 1.

Chapter 5 discusses the findings, limitations, and future work. It interprets the experimental results in the context of real-world deployment, examines the principal limitations of the current system (including the network-isolated collection environment and per-PID tokenization scope), and outlines future research directions including expanded family and variant coverage and production deployment considerations.

\chapter{Background and Literature Review}

This chapter reviews the background and prior work that motivate This thematic paper. It examines modern ransomware behavior, eBPF-based kernel monitoring at the VFS layer, and machine-learning approaches to ransomware detection, including the challenge of generalization across families and deployment environments. It also surveys relevant public datasets and concludes by identifying the research gaps addressed in This thematic paper and positioning the proposed approach within the existing literature.

\section{Ransomware: Characteristics and Behavioral Taxonomy}

This part reviews the main characteristics of ransomware and the behavioral patterns most relevant to detection on Linux.

\subsection{Definition and Threat Landscape}

Ransomware is a category of malware that denies victims access to their own data or systems and demands payment --- typically in cryptocurrency --- for restoration \cite{ref1}. Unlike covertly operating malware such as spyware or banking trojans, ransomware is inherently disruptive: its damage becomes apparent within minutes to hours of execution, making early detection critical. Linux-based cloud and enterprise servers have become particularly prominent targets, as a successful attack on a Linux host can simultaneously disable virtualized workloads, storage services, and containerized applications \cite{ref2}, \cite{ref3}. Families such as Akira \cite{ref4}, BlackBasta \cite{ref5}, IceFire \cite{ref6}, and LockBit \cite{ref7} have each released dedicated Linux encryptors targeting virtualization platforms such as VMware ESXi. The financial impact of ransomware remains substantial, with a growing share of ransomware campaigns specifically targeting Linux infrastructure \cite{ref1}.

\subsection{Conceptual Ransomware Execution Lifecycle}

The conceptual ransomware execution lifecycle on Linux is illustrated in Figure 2.1. It comprises four phases that are common across most modern ransomware families: reconnaissance, key generation, encryption, and extortion.

The exact ordering and implementation of these phases vary across ransomware families \cite{ref2}, \cite{ref8}. During the reconnaissance stage, the malware identifies candidate files and directories for encryption, often prioritizing user documents, databases, and virtual machine images. This preparatory behavior is followed by cryptographic setup, in which the malware initializes the keys or key material required for subsequent encryption; however, the exact mechanism differs across families and may involve locally generated session keys, embedded public-key material, or external key exchange with attacker-controlled infrastructure \cite{ref2}, \cite{ref8}.

The most operationally significant stage for This thematic paper is the encryption phase. At this stage, victim files are read, transformed, and written back either to the original path or to a new file that may subsequently be renamed with a family-specific extension. These operations produce observable VFS-level I/\allowbreak{}O activity, including characteristic sequences of read, write, and rename events, which form the primary behavioral signal exploited by the proposed detection system \cite{ref9}. After encryption, many families enter an extortion stage by dropping ransom notes or otherwise presenting user-facing extortion messages, after which the process either terminates or remains resident \cite{ref2}, \cite{ref8}.

From the perspective of kernel-level behavioral monitoring, these stages are not equally visible. Reconnaissance and extortion may generate relatively sparse or family-specific artifacts, whereas the encryption stage consistently produces the dense file I/\allowbreak{}O patterns most suitable for sequence-based detection. For this reason, the present thesis focuses on modeling the behavioral traces emitted during the encryption stage, while treating the broader lifecycle as contextual background for interpreting ransomware activity on Linux.

\begin{figure}[htbp]
\centering
\includegraphics[width=\linewidth]{fig2-1.png}
\caption{Conceptual ransomware execution lifecycle on Linux.}\label{fig:2-1}
\end{figure}

\subsection{Encryption Strategies and Evasion}

Ransomware families differ substantially in their encryption strategies, and these differences produce distinct, measurable I/\allowbreak{}O signatures. The simplest approach is full-file overwrite, in which the ransomware reads each victim file, encrypts its entire content, and writes the ciphertext back to the same path, producing a symmetric READ--WRITE pair for every file. More evasive families employ selective or partial-file encryption, writing encrypted content only to specific file offsets while leaving other regions intact; this produces lower total write volume and may generate I/\allowbreak{}O patterns indistinguishable from those of legitimate application workloads, causing volume-based and entropy-based detectors to fail \cite{ref8}. A third pattern combines encryption with extension renaming, appending a family-specific suffix after writing, and produces a three-operation READ--WRITE--RENAME sequence per file. Families that perform rapid, high-volume renaming expose a particularly distinctive rename-rate signature. The specific write-size distributions, rename rates, and evasion characteristics of the fourteen families studied in this work are characterized empirically in Chapter 4, based on analysis of raw eBPF VFS traces.

\section{eBPF for Kernel-Level Security Monitoring}

This part reviews eBPF as a kernel-level observability mechanism and considers its relevance to security monitoring for ransomware detection on Linux.

\subsection{Overview of eBPF}

The extended Berkeley Packet Filter (eBPF) executes user-supplied programs within a sandboxed in-kernel virtual machine, providing native-speed instrumentation without the stability risk of loadable kernel modules \cite{ref10}, \cite{ref11}. For security monitoring, the attachment point---rather than eBPF itself---determines the semantic richness and overhead of the resulting event stream. The internal mechanics of bytecode verification, just-in-time compilation, and map-based user-space transfer are well documented in prior work \cite{ref10} and are summarized here only insofar as they constrain sensor design (Section 2.2.4).

\begin{figure}[htbp]
\centering
\includegraphics[width=\linewidth]{fig2-2.png}
\caption{Simplified architecture of the eBPF-based VFS monitoring and ransomware detection pipeline.}\label{fig:2-2}
\end{figure}

\subsection{eBPF Hook Points and Event Sources}

Different kernel attachment points expose different views of file activity. System-call tracepoints capture user-space intent with clean process attribution but can be blinded by asynchronous or buffered I/\allowbreak{}O paths. Block-layer tracepoints expose physical storage operations but lose the link to the originating process, because kernel caching and request merging obscure the file-to-process mapping. The Virtual File System (VFS) layer sits between these extremes, observing logical file operations with full process context while remaining above block-layer caching effects (Figure 2.2).

\subsection{VFS-Level Monitoring for File I/\allowbreak{}O Behavior}

In Linux, the VFS layer provides a uniform interception point for file operations across heterogeneous concrete file systems \cite{ref41}. Attaching eBPF kprobes to vfs\_\allowbreak{}read, vfs\_\allowbreak{}write, and vfs\_\allowbreak{}rename yields per-event metadata---operation type, file path, size, offset, and timestamp---together with the identity of the issuing process. Preserving process-aware event ordering is advantageous for sequence-based detection, because temporal patterns may carry behavioral information that is lost when events are aggregated into summary statistics. This telemetry does not observe writes that bypass the VFS read and write path, such as memory-mapped writeback or io\_\allowbreak{}uring submissions, a coverage boundary examined in Section 5.3.

\subsection{Strengths and Constraints of eBPF Monitoring}

eBPF offers a practical middle ground between user-space monitoring and intrusive kernel modification, combining process-aware context with low overhead \cite{ref10}, \cite{ref11}. Two constraints bear directly on sensor design. First, every program must pass the kernel verifier, which restricts unbounded loops and unsafe memory access and thus limits in-kernel model complexity, motivating the user-space event-filtering and sequence-classification pipeline adopted in Chapter 3. Second, both ransomware and benign high-I/\allowbreak{}O workloads can saturate the kernel-to-user ring buffer, so throughput and event-loss behavior must be measured rather than assumed (Chapter 4).

\section{Machine Learning for Ransomware Detection}

This part reviews the main machine-learning approaches used for ransomware detection and highlights the issues most relevant to This thematic paper.

\subsection{Feature-Based and Storage-Layer Approaches}

Early machine learning approaches to ransomware detection relied on hand-crafted features extracted from file system events, block-level I/\allowbreak{}O logs, or system call traces. Yang and Sahita \cite{ref29} build a behavior-based deep learning detector for crypto-ransomware using I/\allowbreak{}O event patterns and file-content entropy as features, reporting approximately 98\% classification accuracy in controlled experiments but noting that naive models exhibit false-positive rates of approximately 1 -- 3\% in real-world deployment --- motivating adversarial evaluation and more robust feature design. Hirano and Kobayashi \cite{ref26} collect ransomware storage access patterns through a live-forensic hypervisor and train Random Forest, Support Vector Machine, and K-Nearest Neighbor classifiers on five hardware-level I/\allowbreak{}O features, achieving a 98\% F-measure. Gagulic et al. \cite{ref25} similarly explore ML classifiers applied to storage-layer I/\allowbreak{}O traces under the supervision of IBM Research Lab Zurich.

While effective against known families, block-layer and storage-layer approaches share a fundamental limitation: they operate below the file-system abstraction boundary, where file-path resolution is unavailable. Without knowledge of which file is being modified, these approaches cannot filter irrelevant I/\allowbreak{}O or associate encryption activity with a specific process, limiting their applicability to fine-grained behavioral analysis in multi-tenant environments.

A related line of work forgoes trained classifiers in favor of behavioral heuristics computed directly over file operations. Scaife et al. \cite{ref51} introduce CryptoDrop, an early-warning detection system that monitors read and write accesses to protected user directories and accumulates a per-process reputation score from indicators including content dissimilarity (measured via fuzzy hashing), Shannon entropy change, file-type funneling, and file deletion; evaluated against nearly 500 real-world ransomware samples, CryptoDrop detects 100\% of samples with a median loss of only 10 files. Because CryptoDrop reasons over per-process file operations rather than a trained statistical model, it generalizes to previously unseen families without retraining; however, like the storage-layer approaches discussed above, it evaluates each indicator independently at the point of file access rather than modeling the temporal sequence of operations, limiting its ability to distinguish early-stage ransomware behavior from bursty benign I/\allowbreak{}O.

\subsection{Kernel-Level and eBPF-Based Detection}

Kernel-level monitoring provides richer behavioral context by observing I/\allowbreak{}O at the file-system abstraction layer, where both process identity and file path are available. Ahmed et al. \cite{ref14} propose Peeler, which profiles kernel-level events on Windows using the Event Tracing for Windows (ETW) framework. Evaluated against 43 distinct ransomware families --- the largest evaluation dataset at the time of publication --- Peeler achieves 99.52\% detection accuracy with a 0.58\% false-positive rate. Crypto-ransomware is detected within 115.3 milliseconds on average after only one file is encrypted, and transferability tests on 86 samples from 26 previously unseen families yield over 95\% accuracy, demonstrating meaningful cross-family generalization.

On Linux, Brodzik et al. \cite{ref13} implement a decision tree and a multilayer perceptron directly inside eBPF programs, eliminating user-space inference latency. Willers and Philippart \cite{ref17} combine eBPF event capture with ML classification, achieving a 100\% true positive rate across 53 ransomware samples. Hitaj et al. \cite{ref20} propose Minerva, a file-based ransomware detector on Linux engineered for robustness against adversarial evasion, achieving over 99\% detection with 0.52-second average detection latency. van Sloun et al. \cite{ref24} address the practical overhead of kernel-level I/\allowbreak{}O monitoring by proposing a multi-staged detection pipeline that balances detection accuracy against monitoring cost, demonstrating that lightweight first-stage filtering can reduce overhead without sacrificing detection rates. Kostopoulos et al. \cite{ref18} and Kim et al. \cite{ref19} further demonstrate eBPF as a general-purpose security instrumentation platform, applying it to network intrusion detection and container cryptojacking detection respectively.

Beyond kernel- and eBPF-specific systems, sequence-based neural architectures have proven effective for malware detection more broadly when applied directly to ordered event traces rather than aggregate statistics. Luo et al. \cite{ref32} encode Windows API call sequences as behavior graphs and process them with a single-bidirectional graph embedding combined with a Bi-LSTM network trained under a multi-task learning objective, reporting that preserving the temporal ordering of API calls yields higher detection accuracy than four comparison methods across two real-world malware datasets. This finding reinforces the design rationale adopted in This thematic paper: representing VFS-level I/\allowbreak{}O activity as an ordered token sequence, rather than collapsing it into per-process aggregate features, retains the fine-grained temporal signal needed to distinguish ransomware-specific I/\allowbreak{}O patterns from benign activity.

\subsection{Generalization and Deployment Challenges}

A critical dimension of ransomware detection research is generalization to unseen behavior. Reategui et al. \cite{ref12} show that the accuracy of machine-learning-based ransomware detection in block storage varies significantly across storage configurations, file systems, and virtualization environments, and cannot be assumed from same-configuration test sets. Their study concerns cross-configuration generalization; the complementary question of cross-family generalization --- whether a model trained on known families detects unseen ones --- remains largely unaddressed for Linux sequence-based detection.

Kaya et al. \cite{ref15} expose a complementary challenge in general malware detection: models trained on sandbox execution traces achieve approximately 95\% true-positive rate at 1\% false-positive rate on further sandbox traces, but this falls to approximately 17\% on real-world endpoint traces. Such sandbox-to-deployment gaps --- caused by distribution shift, label noise, and environment-specific features --- motivate evaluation methodologies that stress-test generalization beyond same-distribution test sets, in particular cross-family holdout evaluation such as leave-one-family-out (LOFO), in which each ransomware family is withheld from training in turn and used exclusively for evaluation.

Adversarial robustness and concept drift are further deployment concerns \cite{ref30}, \cite{ref31}: adversaries may adapt encryption strategies to evade learned behavioral signatures \cite{ref8}, and the shift toward selective or partial-file encryption (Section 2.1.3) is one such evolution. This thematic paper does not evaluate robustness against deliberate evasion; this is noted as a limitation and a direction for future work.

\begin{table}[htbp]
\centering
\small
\caption{Comparison of selected ransomware and malware detection systems.}\label{tab:2-1}
\begin{tabularx}{\textwidth}{XllXllX}
\hline
\textbf{System} & \textbf{OS} & \textbf{Sensor} & \textbf{ML Model} & \textbf{Stateful} & \textbf{Metrics} & \textbf{Key Limitation} \\
\hline
Hirano \& Kobayashi \cite{ref26} (2019) & Windows & Block I/\allowbreak{}O (hypervisor) & RF /\allowbreak{} SVM /\allowbreak{} KNN & No & F-measure & Block-layer; no file-path resolution \\
Yang \& Sahita \cite{ref29} (2019) & Windows & I/\allowbreak{}O events + entropy & LSTM + SVM & LSTM: Yes; SVM: No & Acc, FPR & High FPR at deployment; Windows only \\
Peeler \cite{ref14} (2021) & Windows & ETW kernel events & Rule-based + ML classifier & No & Acc, FPR & Windows ETW only; no Linux /\allowbreak{} eBPF \\
Gagulic et al. \cite{ref25} (2023) & Any & Storage I/\allowbreak{}O & RF /\allowbreak{} ML classifier & No & Acc & Storage-layer only; Master project \\
Brodzik et al. \cite{ref13} (2024) & Linux & eBPF syscalls & DT /\allowbreak{} MLP & No & Acc, latency & In-kernel ML; no sequence model \\
Willers \& Philippart \cite{ref17} (2024) & Linux & eBPF events & ML classifier & No & TPR & No sequence model; small sample set \\
CLEAR \cite{ref9} (2025) & Windows & NVMe commands & \textmu{}LSTM /\allowbreak{} RF & \textmu{}LSTM: Yes; RF: No & MDR, MBD3 & Windows only; no LOFO evaluation \\
LARM \cite{ref33} (2025) & Linux & Syscall monitoring & ML classifier (multi-feature) & No & TPR, FPR & Fixed-window aggregation; no sequence model \\
Minerva \cite{ref20} (2025) & Linux & File VFS & RF + entropy-based features & No & Acc, FPR & No sequence model; no LOFO \\
RansoGuard \cite{ref21} (2025) & Windows & API call sequences & RNN & Yes & Recall, Acc & Pre-attack API only; Windows only \\
This work (2026) & Linux & eBPF (VFS kprobes) & Stateful \textmu{}LSTM & Yes & MDR, MBD3 & VM-controlled environment; 14 families; single-machine \\
\hline
\end{tabularx}
\end{table}

Note. Acc = accuracy; FPR = false-positive rate; TPR = true-positive rate; MDR = missed detection rate. In This thematic paper, MBD3 is defined as the 75th percentile of cumulative WRITE bytes observed before the first detection of each detected malware process. Metric definitions, sampling units, thresholds, and operating points differ across the cited studies; therefore, the entries summarize the metrics reported by each study and should not be interpreted as direct numerical comparisons.

As shown in Table 2.1, prior work spans a wide range of operating systems, monitoring layers, and model families, from storage-level feature-based classifiers to kernel-level monitoring and recurrent sequence models. However, none of the Linux-based systems included in Table 2.1 explicitly combines VFS-layer monitoring with a stateful sequence model for ransomware detection. This observation is limited to the systems reviewed in this chapter and does not establish the absence of comparable approaches outside the surveyed literature. This positioning motivates the approach adopted in the present thesis, which integrates eBPF-based VFS monitoring with compact event tokenization and stateful sequence classification under a cross-family evaluation setting.

\section{Datasets for Ransomware Research}

Publicly available datasets have played a central role in advancing ransomware detection research by providing reproducible experimental baselines. Table 2.2 summarizes key existing datasets alongside the dataset constructed for This thematic paper, comparing them across dimensions relevant to sequence-based detection on Linux. Among the datasets included in this comparison, none provides the same combination of Linux-based VFS-level event streams, per-event labels, and mixed benign workloads. This observation is limited to the studies and datasets reviewed in this chapter and should not be interpreted as evidence that no comparable dataset exists outside the surveyed literature. These datasets differ primarily in sensor layer and platform (Table 2.2): hypervisor and NVMe captures (RanSAP \cite{ref22}, RanSMAP \cite{ref23}, Ma et al. \cite{ref36}, CLEAR \cite{ref9}) and block- or memory-level traces (DeftPunk \cite{ref34}, Zhu et al. \cite{ref35}) provide low-level I/\allowbreak{}O fidelity but sacrifice file-path resolution, while endpoint-resident monitors (Redemption \cite{ref38}, ShieldFS \cite{ref39}) log filesystem requests on Windows hosts. The dataset constructed for This thematic paper is collected via eBPF VFS kprobes, which yields per-process event streams with full file-path, byte offset, and transfer size at the cost of requiring controlled virtual machine execution environments. Among the surveyed datasets, only Higuchi and Kobayashi \cite{ref37} also adopt eBPF on Linux, although their dataset covers only five ransomware families and does not provide per-event command labels. In Table 2.2, the Public column uses three availability categories: Yes indicates an openly released dataset, Restricted indicates availability subject to stated access conditions, and No indicates that no public release was identified.

\begin{table}[htbp]
\centering
\small
\caption{Comparison of ransomware research datasets.}\label{tab:2-2}
\begin{tabularx}{\textwidth}{lrlXllll}
\hline
\textbf{Dataset} & \textbf{Year} & \textbf{Public} & \textbf{Family /\allowbreak{} variant coverage} & \textbf{OS} & \textbf{Sensor Layer} & \textbf{Event Labels} & \textbf{Mixed Workloads} \\
\hline
This work (2026) & 2026 & No & 14 families /\allowbreak{} 28 binaries & Linux & eBPF VFS (kprobe) & Yes & Yes \\
CLEAR \cite{ref9} & 2025 & Yes & 47 families /\allowbreak{} 137 variants & Windows & NVMe commands & Yes & Yes \\
RanSMAP \cite{ref23} & 2025 & Yes & 6 families & Windows & Block + memory & No & Yes \\
DeftPunk \cite{ref34} & 2024 & No & 13 families & Windows, Linux & Cloud memory & No & Yes \\
Zhu et al. \cite{ref35} & 2024 & No & 10 families & Linux & Storage & No & Yes \\
Ma et al. \cite{ref36} & 2023 & No & 32 families & Windows, Linux & Hypervisor + SSD & No & No \\
Higuchi \& Kobayashi \cite{ref37} & 2023 & No & 5 families & Linux & eBPF & No & No \\
RanSAP \cite{ref22} & 2022 & Yes & 7 families & Windows & Block /\allowbreak{} hypervisor & No & No \\
Redemption \cite{ref38} & 2017 & No & 29 families & Windows & Endpoint & No & No \\
ShieldFS \cite{ref39} & 2016 & Restricted & 11 families & Windows & IRP (filesystem) & No & Yes \\
\hline
\end{tabularx}
\end{table}

Hirano et al. \cite{ref22} introduce RanSAP, a hypervisor-captured dataset of storage-access patterns from seven ransomware families, and extend it to the multimodal RanSMAP \cite{ref23}; both are Windows-based. File-level corpora such as NapierOne \cite{ref28} support entropy-based encrypted-file identification \cite{ref27} but are not behavioral datasets and fall outside the scope of Table 2.2.

Because these traces are collected in virtual machines rather than on production hosts, the dataset shares the sandbox-to-production deployment gap that Kaya et al. \cite{ref15} observe; bridging it through field deployment is beyond the scope of this work. The full collection environment and methodology are detailed in Chapter 3.

\section{Research Gaps and Thesis Positioning}

The preceding review identifies three research gaps that motivate the design and evaluation of This thematic paper.

The first gap concerns the limited integration of Linux eBPF-based VFS-level monitoring with stateful sequence classification. Existing Linux-based ransomware detection systems have demonstrated the practicality of eBPF for runtime monitoring, but they mainly rely on rule-based logic, trap-file monitoring, fixed-window features, or conventional machine-learning models rather than stateful sequence modeling \cite{ref10}, \cite{ref11}, \cite{ref13}, \cite{ref17}. In contrast, storage-oriented ransomware studies provide useful low-level I/\allowbreak{}O evidence, but many operate below the file-system abstraction layer, where file-path context and direct process-level semantics are less available \cite{ref22}, \cite{ref23}, \cite{ref25}, \cite{ref26}. Windows-based systems and CLEAR-style approaches show the value of fine-grained monitoring and stateful modeling \cite{ref9}, \cite{ref14}, but they cannot be directly transferred to Linux because the instrumentation mechanisms and file-system semantics differ across platforms. In Linux, VFS-level monitoring introduces specific technical challenges, including dentry-based path reconstruction, preserving correct process and file-path association at kernel hooks, and handling file-system-specific I/\allowbreak{}O behavior such as ext4 journaling or Btrfs copy-on-write. This thematic paper addresses this gap by combining eBPF VFS-level kprobes with a stateful \textmu{}LSTM classifier, as described in Chapter 3.

The second gap concerns generalization evaluation for sequence-based ransomware detection on Linux kernel-level I/\allowbreak{}O traces. Prior studies have shown that ransomware detection models may not generalize reliably across configurations, environments, or deployment settings \cite{ref12}, \cite{ref15}. However, to the best of the author's knowledge, leave-one-family-out (LOFO) evaluation has not been sufficiently explored for sequence classifiers trained on Linux eBPF VFS-level traces. This is important because a detector that performs well on known ransomware families may still fail when exposed to an unseen family with different encryption behavior. This thematic paper addresses this gap through a fourteen-family LOFO evaluation, presented in Chapter 4.

The third gap concerns the empirical characterization of selective encryption evasion from raw Linux VFS-level I/\allowbreak{}O traces. Hirano and Kobayashi \cite{ref8} demonstrate that adversarial manipulation of encryption parameters can degrade ransomware detection. However, there remains limited empirical characterization of such behavior using Linux VFS-level traces with file-path context and byte-granularity write-size information. This thematic paper addresses this gap by analyzing selective-encryption behavior across fourteen ransomware families and characterizing how such behavior appears in raw VFS-level I/\allowbreak{}O traces, as discussed in Chapter 4.

Together, these gaps define the positioning of This thematic paper. Existing studies demonstrate the value of eBPF, kernel-level monitoring, storage I/\allowbreak{}O analysis, and sequence modeling, but these directions have rarely been combined into a Linux VFS-level ransomware detection pipeline. This thematic paper therefore focuses on collecting process-aware VFS-level I/\allowbreak{}O events with eBPF, representing them as ordered behavioral sequences, and evaluating their effectiveness under family-level generalization and selective-encryption analysis. Chapter 3 presents the system architecture that implements this design, while Chapter 4 presents the experimental evaluation.

\chapter{Methodology}

This chapter presents the methodology used in This thematic paper for the proposed eBPF-based ransomware detection system. It introduces the system architecture, data collection infrastructure, dataset construction process, tokenization scheme, and model architecture. It also describes the LOFO evaluation protocol and the associated metrics, before concluding with a real-time deployment prototype that demonstrates the practical feasibility of the proposed approach.

\section{System Overview}

This part introduces the overall design of the proposed ransomware detection system and the main requirements that shape its architecture.

\subsection{Design Goals and Requirements}

The proposed system is designed to detect file-encrypting ransomware on Linux endpoints in real time. The system addresses five design requirements derived from the operational constraints of production deployment and the limitations identified in existing research \cite{ref9}:

\begin{itemize}
\item Real-time detection: target low-latency streaming detection from VFS events while measuring data exposure before the first Stage 1 detection.
\item Cross-family generalization: the ability to detect previously unseen ransomware families without model retraining, evaluated through leave-one-family-out (LOFO) evaluation.
\item Process-level granularity: classification operates at the individual process level, not at aggregate system level, enabling targeted quarantine of the offending process.
\item No kernel patches required: the system uses eBPF with standard capabilities (CAP\_\allowbreak{}BPF, CAP\_\allowbreak{}PERFMON), avoiding the need for custom kernel modules or patches.
\item Prototype deployability: the system is packaged as a systemd service with a web-based monitoring dashboard and is designed for deployment on Linux hosts with kernel 5.8 or later, the first kernel release to provide the BPF ring buffer (BPF\_\allowbreak{}MAP\_\allowbreak{}TYPE\_\allowbreak{}RINGBUF) and the dedicated CAP\_\allowbreak{}BPF capability on which the sensor relies for low-overhead, least-privilege event delivery. The kernel 5.8 requirement identifies the minimum availability of these facilities, not a guarantee of compatibility across later kernels. The complete pipeline was tested on Ubuntu 24.04 LTS with kernel 6.17.0-14-generic. Kernel symbols, kprobe signatures, BCC behavior, and dentry path reconstruction may vary across releases and distributions; deployment on another kernel therefore requires validation of probe attachment, event-field interpretation, and reconstructed paths.
\end{itemize}

\subsection{System Architecture Overview}

The proposed system implements an end-to-end pipeline for capturing, representing, and classifying process-level file I/\allowbreak{}O behavior in real time. As illustrated in Figure 3.1, the architecture separates kernel-space event capture from userspace processing and consists of five stages.

\begin{itemize}
\item Stage 1 --- Kernel-level event capture: eBPF kprobes attached to vfs\_\allowbreak{}read, vfs\_\allowbreak{}write, and vfs\_\allowbreak{}rename capture file I/\allowbreak{}O events at the Linux Virtual File System (VFS) layer. For each event, the sensor records the operation type, process identifier, process name, byte offset, transfer size, timestamp, and up to seven components of the file path. The resulting event records are transferred from kernel space to userspace through a shared BPF ring buffer.
\item Stage 2 --- Userspace preprocessing and tokenization: A userspace callback reconstructs the file path from the captured dentry components and removes events associated with predefined orchestration and monitoring utilities. The remaining events are grouped by process and converted into pairs of 10-bit tokens. These tokens encode temporal, spatial, operational, and byte-overlap information, as described in Section 3.4.
\item Stage 3 --- Stateful sequence inference: For each monitored process, the token sequence is processed in chunks of 100 events by a stateful \textmu{}LSTM containing 182,849 trainable parameters. The model maintains a separate hidden state for each process and carries this state across consecutive chunks, allowing information from earlier I/\allowbreak{}O activity to influence subsequent predictions without retaining the complete event history.
\item Stage 4 --- Alert decision: The model produces a ransomware probability for each processed chunk. In the research evaluation, a process is considered detected when its probability reaches or exceeds the fixed evaluation threshold of 0.30. The deployment prototype applies a stricter three-layer decision procedure: a probability threshold of 0.85, a three-consecutive-chunk consistency requirement, and a heuristic verification stage based on process and file-path characteristics. The distinction between the evaluation and deployment configurations is detailed in Sections 3.7 and 3.8.
\item Stage 5 --- Response and reporting: When the deployment pipeline confirms an alert, the system records the process and affected-file information for subsequent investigation. The prototype can suspend the suspected process using SIGSTOP and optionally initiate a filesystem snapshot where the underlying storage configuration supports this operation. These response functions demonstrate system integration and are not included in the primary LOFO detection evaluation.
\end{itemize}

The proposed system draws on the sequence-modeling and tokenization concepts introduced by the CLEAR framework \cite{ref9}, which was developed using Windows NVMe command streams collected offline through Xperf. Rather than directly reproducing CLEAR, this work adapts its underlying concepts to Linux by introducing a live eBPF sensor that captures process-level file I/\allowbreak{}O events at the VFS layer. This adaptation addresses the limited availability of Linux-generated ransomware I/\allowbreak{}O traces identified by CLEAR while preserving file-path and process context that is unavailable at the NVMe command layer. Figure 3.1 presents the resulting five-stage architecture and shows the boundary between kernel-space event capture and userspace tokenization, inference, alert confirmation, and response.

\begin{figure}[htbp]
\centering
\includegraphics[width=\linewidth]{fig3-1.png}
\caption{Five-stage end-to-end detection pipeline, from kernel-level eBPF capture to automated response.}\label{fig:3-1}
\end{figure}

\subsection{Design Rationale: Why eBPF at the VFS Layer}

Three candidate instrumentation layers were considered for kernel-level I/\allowbreak{}O monitoring: the system call interface, the Virtual File System (VFS) layer, and the block device layer. Table 3.1 summarizes their trade-offs in terms of path visibility, PID attribution, and expected monitoring overhead.

\begin{table}[htbp]
\centering
\small
\caption{Comparison of kernel I/\allowbreak{}O monitoring layers.}\label{tab:3-1}
\begin{tabularx}{\textwidth}{lllX}
\hline
\textbf{Layer} & \textbf{Path Visibility} & \textbf{PID Attribution} & \textbf{Expected relative overhead} \\
\hline
Syscall & Full path & Direct & Expected higher; depends on the traced syscall set (not measured in this study) \\
VFS (selected) & Full path via dentry & Direct & Expected moderate for the selected VFS hooks (not measured in this study) \\
Block device & None (sector only) & Requires BPF map join & Expected lower at the capture point; attribution joins add cost (not measured in this study) \\
\hline
\end{tabularx}
\end{table}

The VFS layer was selected because it provides both full file path visibility (via dentry traversal) and direct process attribution (via bpf\_\allowbreak{}get\_\allowbreak{}current\_\allowbreak{}pid\_\allowbreak{}tgid), while incurring moderate overhead by firing probes only on file operations rather than on every system call. The block device layer, while lower overhead, lacks path and PID information without additional BPF map infrastructure to join inode-to-path mappings a complexity CLEAR notes is costly even for offline processing \cite{ref9}.

WaR and RaR byte-overlap features are implemented offline rather than inside the eBPF probe itself. Computing these features in-kernel would require tracking the full history of prior reads on every file through BPF maps, which the CLEAR paper characterizes as prohibitively expensive in its CLEAR-NVMe section \cite{ref9}. Instead, WaR and RaR are computed at tokenization time from the captured event stream using CLEAR's same-file byte-range overlap criterion, and encoded in the Token2 aux bits as defined in section 3.4.1.

\subsection{Threat Model and Assumptions}

The proposed system targets Linux file-encrypting ransomware that generates observable I/\allowbreak{}O patterns at the Virtual File System (VFS) layer \cite{ref50}. The threat model assumes that the ransomware performs a sequence of READ (to obtain plaintext), WRITE (to overwrite with ciphertext), and optionally RENAME (to change the file extension) operations on victim files, and that these operations are issued through the standard kernel entry points instrumented by the eBPF sensor --- specifically vfs\_\allowbreak{}read, vfs\_\allowbreak{}write, and vfs\_\allowbreak{}rename. The threat model further assumes that the ransomware process is not privileged to disable or unload the eBPF instrumentation, and that the operating system kernel and the eBPF subsystem are intact at detection time. The reported false-positive rates assume benign workloads drawn from the categories represented in the training set (archiving, legitimate encryption, secure deletion, developer activity, document editing, file copying, and simulated web browsing); operating-environment workloads outside this distribution may exhibit different false-positive characteristics and would require workload-specific heuristic tuning.

The following are explicitly out of scope. First, kernel-level rootkits that disable, unload, or subvert eBPF probes; fileless malware operating entirely in memory; pure data exfiltration without on-disk encryption; and ransom-note droppers that do not encrypt files. Second, I/\allowbreak{}O paths that bypass the instrumented kernel entry points: vectored I/\allowbreak{}O issued through pwritev and preadv, asynchronous submission through io\_\allowbreak{}uring, memory-mapped writeback via mmap and msync, and direct-mode writes via O\_\allowbreak{}DIRECT. The current sensor does not observe these paths, and Guunra and a subset of Akira variants empirically exhibit this evasion behavior (see Section 4.5.2). Third, adversarial behaviors within userspace that exploit the heuristic gate, including process-name spoofing via prctl(PR\_\allowbreak{}SET\_\allowbreak{}NAME) to match the trusted-process allowlist, encrypted-extension obfuscation that avoids the eighteen-entry ransom-extension list, target-path mimicry that emits the majority of writes into system directories before pivoting to user data, and intermittent or low-and-slow encryption strategies that emit writes too sparsely for sustained-chunk detection. Heuristic robustness against these adversarial moves is not evaluated. Fourth, sub-quantization writes: an attacker that writes fewer bytes per operation than the smallest size bucket can pin every WRITE token to the same size quantum, defeating the size-based behavioral signal. Extending the sensor hook set is identified as future work; the effect of token-resolution limits is examined through the SougoLock case in Section 4.5.3 and Appendix F.6 and is summarized in Section 5.1.1.

\section{Data Collection}

This section describes the infrastructure used to collect the dataset of ransomware and benign file-I/\allowbreak{}O events. The environment comprises an isolated virtual machine, an eBPF-based kernel sensor, a scripted benign workload generator, an automated orchestration pipeline, and a CSV-to-Parquet conversion stage.

\subsection{Isolated Virtual Machine Environment}

All data collection was performed in an isolated virtual machine to prevent ransomware escape to the host system. The VM was configured as follows: Oracle VirtualBox hypervisor, Ubuntu 24.04 LTS guest with kernel 6.17.0-14-generic, 14 vCPUs (allocated from the host's available cores), 8 GB RAM, and NAT networking. A secondary 512 GB NVMe SSD was mounted as a dedicated victim file storage disk. In this context, isolated refers to containment of file-system effects within the VM and restoration from a clean snapshot; it does not mean network air-gapping. NAT networking remained enabled for SSH-based orchestration and SFTP result retrieval. Network propagation and unrestricted outbound communication were not evaluated and remain outside the threat model.

Victim files were organized into five sets (Set 1 - 5), each containing files from three sources: the NapierOne standard ransomware research corpus \cite{ref28}, images from the DiffusionDB dataset \cite{ref42}, and Python-generated synthetic files with controlled entropy profiles. Data collection was conducted between March and April 2026 using the isolated VM environment described above. Per-run metadata (source family, hash, victim-file set, event counts, duration) is preserved in the project repository. Table 3.2 summarizes the victim file sets.

\begin{table}[htbp]
\centering
\caption{Victim file sets used for data collection.}\label{tab:3-2}
\begin{adjustbox}{max width=\textwidth}
\begin{tabular}{rrrll}
\hline
\textbf{Set} & \textbf{Size (GB)} & \textbf{Files} & \textbf{Avg File Size} & \textbf{Sources} \\
\hline
1 & 29 & 20,129 & 1.4 MB & NapierOne + Synthetic \\
2 & 172 & 117,139 & 1.5 MB & NapierOne + DiffusionDB \\
3 & 164 & 155,101 & 1.1 MB & NapierOne + DiffusionDB \\
4 & 73 & 170,170 & 430 KB & DiffusionDB \\
5 & 30 & 70,314 & 430 KB & DiffusionDB + Synthetic \\
\hline
\end{tabular}
\end{adjustbox}
\end{table}

Before each data collection run, the VM was restored to a clean snapshot using VBoxManage, ensuring that (a) no residual encrypted files contaminated subsequent runs, and (b) the system state was identical across all experiments.

\subsection{eBPF Sensor Design}

The data collection sensor uses the BCC (BPF Compiler Collection) framework \cite{ref41} to attach three kernel probes (kprobes) at the VFS layer: vfs\_\allowbreak{}read (opcode = 0), vfs\_\allowbreak{}write (opcode = 1), and vfs\_\allowbreak{}rename (opcode = 5). Each probe records the file offset, size, path, and process identity required for downstream tokenization.

Each captured event is first staged in a one-record per-CPU array (BPF\_\allowbreak{}PERCPU\_\allowbreak{}ARRAY) to avoid exceeding the eBPF stack limit and is then submitted to userspace through BPF\_\allowbreak{}PERF\_\allowbreak{}OUTPUT. The userspace reader opens a perf buffer with page\_\allowbreak{}cnt = 8192 pages per CPU and registers a lost-event callback. Full path reconstruction walks up to seven dentry levels to balance verifier safety and path fidelity. The real-time deployment prototype described in Section 3.8 uses a shared BPF ring buffer instead.

The sensor includes two measures that impede casual interference: the sensor binary and its parent directory are marked immutable via chattr +i, and the sensor runs under a dedicated systemd unit. Neither measure is tamper-proof against an adversary with root privileges, who can clear the immutable attribute with chattr -i or terminate the systemd unit directly; consistent with the threat model of Section 3.1.4, the sensor assumes a non-root adversary or one that does not actively target the monitoring infrastructure, and the system does not claim a formal privilege-separation boundary between the sensor and a privileged attacker. Implementation details --- including the full opcode map, buffer sizing rationale, dentry-walk algorithm, and kprobe attach sequence --- are given in Appendix C.1.

\subsection{Benign Workload Generator}

To ensure that the benign class contains realistic file-I/\allowbreak{}O patterns --- including archiving, legitimate user-initiated encryption, and developer workflows that naturally produce Write-after-Read sequences --- a Python-based benign workload generator is used both alongside the malware in mixed-mode captures and independently in benign-only captures.

The generator draws files from the configured file pools described in Appendix B and selects uniformly among seven task methods: archiving, legitimate encryption, secure deletion, developer activity, document editing, disk operations, and simulated web browsing. Appendix B documents the file-pool mapping, implementation details, scheduling behavior, and reproducibility limitations.

\subsection{Automated Orchestration Pipeline}

The data collection process was automated using a Python-based orchestrator that manages the complete lifecycle of each experiment run through six phases: (1) environment preparation via VBoxManage snapshot restore and SSH bootstrap, (2) sensor launch, (3) malware execution from an encrypted archive, (4) a 600-second collection loop with early termination, (5) graceful SIGINT cleanup and SFTP retrieval of the raw CSV, and (6) a quality gate that voids runs with fewer than 50 malware events.

The orchestrator deploys malware samples and an optional benign workload generator (Section 3.2.3) into an isolated Ubuntu 24.04 LTS virtual machine via SSH. Figure 3.2 depicts the data flow. Phase-level timing budgets, retry policies, and the full command sequence for each phase are documented in Appendix D.3.

\begin{figure}[htbp]
\centering
\includegraphics[width=\linewidth]{fig3-2.png}
\caption{Data collection components showing orchestrator, VM, eBPF sensor, ring buffer, event logger, and Parquet output stages.}\label{fig:3-2}
\end{figure}

\subsection{CSV-to-Parquet Conversion and Labeling}

The raw CSV sensor output is converted to Apache Parquet format through a multi-step pipeline that parses with Pandas (C-engine, Python fallback for malformed rows), applies event-level labeling via process-name matching, and emits three output artifacts per run: recording.parquet, process\_\allowbreak{}tree.parquet, and recording\_\allowbreak{}metadata.json. The detailed labeling methodology is described in Section 3.3.3; the Pandas parsing configuration, malformed-row recovery, and full schema of the three artifacts are given in Appendix C.3.

\section{Dataset}

This section presents the dataset used for training and evaluation. It covers the 14 ransomware families included, the overall dataset composition, the process-level labeling methodology, data quality assurance, class imbalance handling, and the noise-filtering procedure applied to malware recordings.

\subsection{Ransomware Families}

The dataset comprises 14 Linux ransomware families representing diverse encryption strategies and I/\allowbreak{}O behavioral patterns.

The evaluation set was assembled from a broader candidate pool of 63 Linux ELF ransomware binaries collected from publicly tracked threat-intelligence repositories and grouped into 30 candidate families. Each binary was executed inside an isolated sandbox virtual machine; candidates that failed to detonate --- typically due to missing command-line flags, hard-coded affiliate passwords, or platform-specific dependencies, such as ESXi runtime --- were excluded from training and evaluation. After this triage, 14 families comprising 28 binaries produced 494 productive sandbox runs and constitute the final dataset reported in This thematic paper. Excluded families include Mirai (network botnet, not ransomware), BrainCipher (ESXi-only encryptor that produced no encryption activity on the Linux VM used in this study), GonnaCry and ItsSoEasy (legacy samples without productive runs), Qilin (requires an affiliate password not available to this study), and Trigona (binary failed to detonate under standard invocation). A complete per-binary inclusion table --- listing SHA-256 hash, architecture, status, and exclusion reason --- is provided in Appendix A.

\begin{table}[htbp]
\centering
\small
\caption{Ransomware family characteristics.}\label{tab:3-3}
\begin{tabularx}{\textwidth}{lrllrX}
\hline
\textbf{Family} & \textbf{Runs} & \textbf{Malware-Labeled Events (Million)} & \textbf{RENAME \%} & \textbf{Median Write (Byte)} & \textbf{Pattern} \\
\hline
Akira & 137 & 128.8 & 8.4\% & 9 & Bimodal (log-spam + encrypt, 11 variants) \\
Bert & 20 & 27.2 & 7.2\% & 3,174 & Clean bulk encryption \\
BlackBasta & 39 & 232.8 & 0.0\% & 64 & Small fixed-chunk, no rename \\
Guunra & 29 & 7.5 & 100.0\% & 0 & RENAME-only (no writes captured) \\
HelloKitty & 38 & 13.0 & 9.2\% & 240 & Medium-chunk encryption \\
IceFire & 35 & 32.4 & 8.6\% & 520 & Medium-chunk encryption \\
LockBit & 39 & 42.6 & 14.2\% & 324 & Variable-size (P95 = 2.3 MB) \\
Monti & 22 & 4.2 & 22.9\% & 35,649 & Clean bulk encryption \\
Nitrogen & 35 & 47.7 & 5.9\% & 794 & Medium-chunk encryption \\
RAWorld & 10 & 5.3 & 20.0\% & 32 & Small-chunk in-place \\
RansomEXX & 38 & 16.0 & 15.5\% & 1,616 & Clean bulk encryption \\
Royal & 20 & 17.4 & 12.1\% & 512 & Medium-chunk encryption \\
Sodinokibi & 22 & 23.4 & 6.8\% & 3,191 & Clean bulk encryption \\
SougoLock & 10 & 0.4 & 0.0\% & 12,544 & Clean bulk encryption, low-volume \\
TOTAL & 494 & 598.6 &  &  &  \\
\hline
\end{tabularx}
\end{table}

Note. The event values report rows with Label = 1 in the selected malware-session Parquet recordings. They exclude Label = 0 events and the 188 standalone benign-workload recordings. The total is 598,613,562 events, rounded to 598.6 million in the table.

The 14 families fall into three behavioral groups based on their observed per-run encryption patterns and WaR byte-overlap frequencies, as characterized in Table 3.3 and Table 4.5. The encryption-dominant group (G1: Bert, BlackBasta, HelloKitty, IceFire, Nitrogen, Sodinokibi, SougoLock) is characterized by sustained file-encryption activity, from small fixed-chunk in-place overwrites to bulk and medium-chunk encryption; the prototype must separate high-volume write bursts from benign editor or compiler I/\allowbreak{}O. The rename-dominant group (G2: Guunra) relies almost exclusively on a temporary-file-then-rename strategy, producing a RENAME-heavy event sequence with negligible ciphertext bytes written to the target path before detection. The bimodal group (G3: Akira, LockBit, Monti, RAWorld, RansomEXX, Royal) interleaves both patterns within a single run or exhibits varying strategies across runs; streaming evaluation captures intermittent high-probability windows between encryption bursts. These groups are carried through the per-family results in Chapter 4 so that performance can be read against behavioral rather than alphabetical ordering. The groupings are manually assigned based on observed VFS signal profiles and are not the output of a formal clustering algorithm.

\subsection{Dataset Composition}

The final dataset comprises 682 usable Parquet recordings: 494 malware-session recordings collected across fourteen ransomware families and 188 standalone benign-workload recordings. Direct counting of rows with Label = 1 in the selected malware-session Parquet recordings yielded 598,613,562 malware-labeled VFS I/\allowbreak{}O events (approximately 598.6 million). This count excludes Label = 0 events and the 188 standalone benign-workload recordings; therefore, it should not be interpreted as the total number of event rows across all 682 recordings.

The 494 malware runs represent 28 unique SHA-256 binary samples across 14 families, with per-family run counts ranging from 10 (RAWorld, SougoLock) to 137 (Akira). Akira's elevated run count reflects the fact that 11 distinct hash variants were collected, each exercised across all five victim sets in both modes. The per-family run distribution is summarized in Table 3.3; the provenance of this spread and the rationale for not rebalancing it are discussed in Section 4.2.

Each malware run directory contains two Parquet files with the same execution provenance but different event composition. The recording.parquet file preserves the events retained from the sensor output, whereas recording\_\allowbreak{}filtered.parquet removes events associated with the orchestration and monitoring utilities described in Section 3.3.6. Training and evaluation use one selected malware recording per run, preferring recording\_\allowbreak{}filtered.parquet when available, together with the standalone benign recordings. The common event schema contains Timestamp, Process ID, Process Name, OpCode, Path Name, Offset (Bytes), Size (Bytes), and Label. The WaR, RaR, and RENAME auxiliary features are computed during tokenization and are not stored as authoritative columns in the selected Parquet recordings.

\section{.3  Labeling Methodology}

Event-level ground-truth labels were assigned within the controlled malware sandbox. For each malware-session recording, two lowercase labeling keywords were obtained from the run metadata: the first twelve characters of the sample's SHA-256 hash and the ransomware family name. The lowercase Process Name value of each captured event was tested for the presence of either keyword using case-insensitive substring matching. An event was assigned Label = 1 when its process name contained at least one of these run-specific identifiers; otherwise, it was assigned Label = 0. The rule was applied independently to every event rather than propagating a label to every event associated with the same PID.

Event-level name matching replaced an earlier PID-tree propagation method that incorrectly assigned malware labels to system utilities, including dpkg, snap, uname, and dirname, when those utilities appeared within the launcher's process tree. The revised rule prevents such labels from being inherited solely through PID or parent--child relationships. Figure 3.3 summarizes the resulting workflow from controlled execution and event capture to process-name-based labeling, noise removal, quality auditing, and generation of the labeled dataset.

The labeling rule was applied to all 494 malware-session recordings. The 188 standalone benign-workload recordings were assigned Label = 0 by construction and did not undergo malware-name matching. The quality audit described in Section 3.3.4 checked that events assigned Label = 1 contained the expected family name or run-specific hash prefix. However, no formal adversarial process-name-spoofing test or exhaustive collision analysis was performed.

Linux process names are length-limited and can be modified by a running process. Substring matching can also produce a false label if an unrelated process name happens to contain a family keyword or hash prefix. These risks are reduced in the controlled collection environment because the executed sample and its SHA-256 hash are known for each run, but they are not eliminated. Accordingly, process-name matching is used only to construct experimental ground-truth labels; it is not presented as a production mechanism for identifying ransomware, and the process name is not included in the \textmu{}LSTM input tokens.

HelloKitty illustrates an edge case involving process execution state. Each missed launcher PID initially appeared under the nohup process name and subsequently executed the HelloKitty binary, after which the recorded process name contained the sample's hash prefix. Under the event-level rule, the initial nohup events were assigned Label = 0, whereas the events recorded after the name changed to the malware hash were assigned Label = 1. Each of the twenty short-lived launcher processes contained exactly sixteen malware-labeled events and no writes to the victim-file corpus. The corresponding main encryption processes were retained as malware and were detected in all thirty-eight HelloKitty runs.

The primary per-process MDR therefore remains 34.48\% because these launcher processes are counted as missed malware processes. A supplementary encrypter-only analysis applies a family-independent exclusion criterion: a malware-labeled process is excluded only when it contains fewer than seventeen malware-labeled events and performs no writes to the victim-file corpus. This criterion excludes the twenty HelloKitty launcher processes while retaining processes that exhibit victim-file encryption activity, yielding an encrypter-only MDR of 0.00\%. Because the threshold was selected after observing the deterministic sixteen-event launcher pattern, this result is reported as a post-hoc diagnostic sensitivity analysis rather than as the primary performance estimate.

\begin{figure}[htbp]
\centering
\includegraphics[width=\linewidth]{fig3-3.png}
\caption{Data-labeling pipeline from VFS event collection to the audited labeled dataset.}\label{fig:3-3}
\end{figure}

Note. Stage (c) applies case-insensitive substring matching between each event's process name and the run-specific SHA-256 prefix or family keyword. The rule is used for offline ground-truth construction and is not part of \textmu{}LSTM inference.

\subsection{Data Quality Audit}

A data quality audit verified all 494 run directories against five criteria: schema consistency, label correctness, PID stability, temporal coverage, and event-count floor. A summary of the audit is presented here, while the full criteria, per-criterion pass/\allowbreak{}fail thresholds, and the two quarantined runs are tabulated in Appendix D.1.

\subsection{Class Imbalance}

The dataset exhibits severe class imbalance at the process level. Established approaches for learning from imbalanced data motivate the weighted-loss and subsampling strategies applied in This thematic paper \cite{ref49}. Because training operates on per-process examples while the model consumes per-event token streams, process-level and event-level distributions are treated as distinct analytical scopes.

At the process level, any Process ID containing at least one Label = 1 event is classified as a malware process; every other PID is treated as benign. Under this definition, the raw process-level ratio is approximately 2,518:1 (benign:malware) across the training partition (see Table 3.4).

\begin{table}[htbp]
\centering
\small
\caption{Process counts by source partition.}\label{tab:3-4}
\begin{tabularx}{\textwidth}{Xrrr}
\hline
\textbf{Source partition} & \textbf{Processes} & \textbf{Benign} & \textbf{Malware} \\
\hline
Malware Ransomware-only runs & 220,317 & 220,040 & 277 \\
Malware Mixed runs & 567,859 & 567,621 & 238 \\
Benign-only runs & 509,091 & 509,091 & 0 \\
Total & 1,297,267 & 1,296,752 & 515 \\
\hline
\end{tabularx}
\end{table}

Note. Process counts are reported before fold-specific benign subsampling and minimum-sequence filtering. Event totals are not included because the raw-capture, label-aware, and post-filter counts represent different analytical scopes and are reported separately in Section 3.3.2.

At the process level, the complete pre-filter pool contains 1,296,752 benign processes and 515 malware processes, corresponding to a benign-to-malware ratio of approximately 2,518:1. This raw process-level ratio does not represent the effective class ratio presented to the model in each fold. As described in Section 3.6.2, fold-specific benign subsampling, minimum-sequence filtering, and weighted loss are applied during training. Event-level and process-level class distributions are therefore reported separately and should not be interpreted as interchangeable measures of class imbalance.

Three complementary strategies are applied: 
(1) Weighted loss is applied through BCEWithLogitsLoss using a bounded positive-class weight, pos\_\allowbreak{}weight = min(max(1.0, $\rho$/\allowbreak{}2), 10), where $\rho$ denotes the benign-to-malware process ratio in each fold's training partition.

(2) Benign sub-sampling: the MAX\_\allowbreak{}BENIGN\_\allowbreak{}PROCS = 5000 cap bounds benign processes per fold, giving a mean post-cap ratio of approximately 10.5:1 (benign:malware), reduced to about 7.5:1 after the minimum-sequence filter of 10 tokens, equivalent to five accepted VFS events.

(3) Recording-level partitioning: all processes extracted from the same Parquet recording are retained in the same partition. In each LOFO fold, all recordings from the held-out family are assigned exclusively to the test set. Recordings from each of the remaining thirteen families are split independently into training and validation subsets, while benign recordings are randomly divided among the training, validation, and test subsets using a fixed seed. The partitions are recording-disjoint but are not grouped by binary hash, victim-file set, or acquisition provenance.

\subsection{Noise Filtering}

A secondary filtering layer removes events produced by orchestrator utility processes (wc, pgrep, ps, awk, top, dirname) that are invoked to manage malware execution and monitoring. These short-lived system utilities contribute heavy-tail I/\allowbreak{}O that is unrelated to the ransomware behavior under study. The full process-name blacklist, its justification, and the corresponding before/\allowbreak{}after WaR/\allowbreak{}RaR rates are provided in Appendix C.2.

\section{Tokenization Scheme}

This part defines the tokenization scheme that converts raw VFS I/\allowbreak{}O events into fixed-width integer tokens suitable for \textmu{}LSTM input. It first presents the 10-bit token layout and its field definitions, and then examines the effective signal carried by each field.

\subsection{10-bit Token Layout}

Each I/\allowbreak{}O event is encoded into two 10-bit tokens (VOCAB\_\allowbreak{}SIZE = 1,024), following a tokenization scheme adapted from the CLEAR framework \cite{ref9}. Token1 encodes temporal, size, and opcode information; Token2 encodes spatial position and behavioral auxiliary features. Table 3.5 defines each field.

\begin{table}[htbp]
\centering
\small
\caption{10-bit token field definitions.}\label{tab:3-5}
\begin{tabularx}{\textwidth}{llrX}
\hline
\textbf{Token} & \textbf{Field} & \textbf{Bits} & \textbf{Encoding} \\
\hline
Token1 & index & 1 & Always 0 (identifies Token1) \\
 & delta\_\allowbreak{}t & 4 & floor(log2(ms + 1)), clamped [0, 15] \\
 & size\_\allowbreak{}q & 4 & floor(log2(bytes + 1)), clamped [0, 15] \\
 & opcode\_\allowbreak{}q & 1 & min(opcode, 1): READ=0, WRITE/\allowbreak{}RENAME=1 \\
Token2 & index & 1 & Always 1 (identifies Token2) \\
 & offset\_\allowbreak{}msb & 4 & (offset /\allowbreak{} 4096) \textgreater{}\textgreater{} 2, masked to 4 bits \\
 & offset\_\allowbreak{}lsb & 2 & (offset /\allowbreak{} 4096) \& 0x3 \\
 & aux & 3 & bit2=RENAME, bit1=WaR, bit0=RaR \\
\hline
\end{tabularx}
\end{table}

RENAME events share the same opcode\_\allowbreak{}q = 1 as WRITE events; the distinguishing signal is instead carried by aux bit2 (rename\_\allowbreak{}flag) in Token2. READ events set opcode\_\allowbreak{}q = 0 and do not activate any auxiliary bit.

The log-quantized fields (delta\_\allowbreak{}t and size\_\allowbreak{}q) compress the wide dynamic range of I/\allowbreak{}O timing and sizes into 4-bit values. The file offset is first converted to 4,096-byte page units, then split into two sub-fields using standard bit manipulation: the upper 4 bits (offset\_\allowbreak{}msb = pages \textgreater{}\textgreater{} 2) capture coarse file position, while the lower 2 bits (offset\_\allowbreak{}lsb = pages \& 0x3) capture fine-grained page alignment. Together, the 6-bit offset field can represent 64 distinct page regions within a file. The auxiliary bits encode byte-range overlap flags following CLEAR's OV\_\allowbreak{}WaR criterion \cite{ref9}: WaR (Write-after-Read) is set to 1 when a WRITE operation's byte range [offset, offset+size] overlaps with the most recent READ on the same file; RaR (Read-after-Read) is set analogously for consecutive READ operations on the same file. The overlap is tracked per-PID using a last-read map keyed by file path, with no cross-process aggregation. These flags provide behavioral signals characteristic of encrypt-in-place ransomware patterns, where the malware reads a file block, encrypts it, and writes the ciphertext back to the same byte range.

The complete token layout is shown in Figure 3.4: Token 1 carries the temporal delta, I/\allowbreak{}O size, opcode, and a token-index bit; Token 2 carries the MSB and LSB offset fields together with the three auxiliary flags (RENAME, WaR, RaR). The figure also presents a worked example for a 12,288-byte WRITE at offset 8,192 occurring 50 ms after the previous event, showing the bit values assigned to each field.

\begin{figure}[htbp]
\centering
\includegraphics[width=\linewidth]{fig3-4.png}
\caption{Layout of the 10-bit token pair.}\label{fig:3-4}
\end{figure}

Three deviations from the CLEAR specification are introduced at the point of tokenization and should be made explicit. First, the prototype uses six offset bits at 4 KiB granularity, covering 256 KiB before the offset wraps, whereas CLEAR uses six offset bits at 2 MiB granularity, covering 128 MiB before the offset wraps. The empirical consequence of this offset-encoding difference on the SougoLock fold is documented in Section 4.5.3 and Table F.5. Second, the third auxiliary bit in the second token encodes a categorical RENAME flag rather than the byte-overlap Read-After-Write indicator used by CLEAR; the bit is set when the opcode is vfs\_\allowbreak{}rename and is otherwise zero. This deliberate substitution reflects the prominence of rename operations in the observed Linux ransomware behaviors and is not a CLEAR-faithful replication. Third, the Write-after-Read auxiliary bit tracks only the latest read interval per path, whereas the CLEAR specification requires consume-on-write semantics over the complete read-interval history; the empirical effect of this difference is examined in Section F.6. These deviations are stated explicitly so that the comparison with CLEAR is interpreted as an inspired adaptation to Linux VFS events rather than as a faithful reimplementation.

\subsection{Effective Signal Analysis}

The tokenizer remaps sensor opcodes to the model encoding (READ = 0, WRITE = 1) before quantization. Each accepted event contributes 18 event-feature bits distributed across two 10-bit token IDs, with one index bit in each token distinguishing Token 1 from Token 2. The event-feature fields are delta\_\allowbreak{}t (4 bits, I/\allowbreak{}O velocity), size\_\allowbreak{}q (4 bits, write-size patterns), opcode\_\allowbreak{}q (1 bit, read versus write), offset (6 bits, spatial file position), RENAME (1 bit, rename behavior), WaR (1 bit, write-after-read overlap), and RaR (1 bit, read-after-read overlap). In CLEAR's block-layer setting, the ablation study reported offset and delta\_\allowbreak{}t as the strongest detection signals, with auxiliary overlap features contributing a smaller but non-negligible effect. This result is cited as directional evidence only, as the VFS-layer setting differs from CLEAR's block-layer measurements.

\section{Model Architecture}

This part describes the neural network architecture and its training configuration. It first presents the stateful \textmu{}LSTM design, then explains the truncated backpropagation through time (TBPTT) training procedure with per-process gradient accumulation, and finally lists the hyperparameters used in all experiments.

\subsection{Stateful \textmu{}LSTM}

The detection model is a stateful \textmu{}LSTM comprising five layers: an embedding layer (1,024 $\times$ 128), a single-layer LSTM (128 $\times$ 64) \cite{ref43}, a dropout layer (p = 0.3) \cite{ref44}, and two fully connected layers (64 $\times$ 32 with ReLU activation, followed by 32 $\times$ 1). The total parameter count is 182,849, making the model suitable for CPU-based inference on resource-constrained endpoints. The model processes one process at a time (batch\_\allowbreak{}size = 1), with the LSTM hidden state (h, c) carried across sequential chunks to maintain temporal context over the entire process lifetime.

\subsection{Truncated Backpropagation Through Time}

Training employs Truncated Backpropagation Through Time (TBPTT) \cite{ref45} with a chunk size of 100 events (200 tokens). Per-process gradient accumulation divides the loss by the number of chunks before calling optimizer.step(), ensuring that every process contributes equal gradient magnitude regardless of its I/\allowbreak{}O length. The hidden state is detached between chunks (hidden = (h.detach(), c.detach())) to prevent gradient flow across chunk boundaries. All training is performed in float32 precision.

\subsection{Hyperparameters}

\begin{table}[htbp]
\centering
\small
\caption{Model hyperparameters.}\label{tab:3-6}
\begin{tabularx}{\textwidth}{lXX}
\hline
\textbf{Parameter} & \textbf{Value} & \textbf{Rationale} \\
\hline
VOCAB\_\allowbreak{}SIZE & 1,024 & 2\textasciicircum{}10 (10-bit tokens) \\
EMBED\_\allowbreak{}DIM & 128 & CLEAR default \\
HIDDEN\_\allowbreak{}SIZE & 64 & CLEAR default \\
NUM\_\allowbreak{}LAYERS & 1 & Single-layer LSTM; additional layers did not improve validation loss \\
CHUNK\_\allowbreak{}SIZE & 200 tokens (100 events) & available GPU memory budget \\
LEARNING\_\allowbreak{}RATE & 0.001 & Adam default \\
DROPOUT & 0.3 & Regularisation \\
GRAD\_\allowbreak{}CLIP & 1.0 & Prevent exploding gradients \\
EPOCHS & 10 (early stop) & Patience = 3 \\
pos\_\allowbreak{}weight & min(max(1.0, ratio/\allowbreak{}2.0), 10.0) & Capped to prevent FP explosion \\
BATCH\_\allowbreak{}SIZE & 1 & Stateful \textmu{}LSTM (per-process) \\
MAX\_\allowbreak{}BENIGN\_\allowbreak{}PROCS & 5,000 & available GPU memory budget \\
MIN\_\allowbreak{}TOKENS\_\allowbreak{}THRESHOLD & 10 tokens & Equivalent to five accepted VFS events; excludes shorter processes \\
SEED & 42 & Reproducibility \\
EVAL\_\allowbreak{}THRESHOLD & 0.30 & Streaming evaluation (Section 3.7.1): maximize recall \\
DEPLOY\_\allowbreak{}THRESHOLD & 0.85 & Real-time Layer 1 (Section 3.8.2): combined with Layer 2 and Layer 3 heuristics \\
\hline
\end{tabularx}
\end{table}

\section{Training Methodology}

This part describes the training methodology used in This thematic paper, including the leave-one-family-out (LOFO) protocol, the three-way data split, the loss function and class-weighting strategy, and the early stopping and learning-rate schedule.

\subsection{Leave-One-Family-Out Protocol}

The leave-one-family-out (LOFO) protocol produced fourteen folds, with one ransomware family withheld as the test family in each fold. All recording files associated with the held-out family were excluded from training and validation. Consequently, no recording from the held-out family was used for model fitting, hyperparameter selection, or early stopping.

The remaining thirteen ransomware families were divided into training and validation partitions at the recording-file level. The split was performed separately within each family using a fixed random seed of 42, with 80\% of the recording files assigned to training and 20\% assigned to validation. Benign recording files were independently divided into training, validation, and test partitions using a 70/\allowbreak{}10/\allowbreak{}20 random split with the same seed. Figure 3.5 illustrates this procedure using the BlackBasta fold as an example.

\begin{figure}[htbp]
\centering
\includegraphics[width=\linewidth]{fig3-5.png}
\caption{Leave-One-Family-Out (LOFO) evaluation protocol.}\label{fig:3-5}
\end{figure}

The malware training-validation split was not grouped by SHA-256 binary hash, victim-file set, or collection provenance. Therefore, recordings generated by the same binary may occur in both the training and validation partitions. The validation results should consequently be interpreted as within-family and potentially within-variant model-selection measurements. This limitation does not place recordings from the held-out family in the training or validation partitions, but it may make validation performance more optimistic than a hash-grouped validation design.

\subsection{3-Way Data Split}

A three-way partition was used to separate model training, model selection, and held-out-family evaluation. The split unit was an individual Parquet recording rather than an individual process, event, binary hash, or victim file. All processes extracted from the same recording therefore remained in the same partition, preventing process-level leakage between training, validation, and test sets.

For each LOFO fold, all recordings belonging to the held-out ransomware family were assigned exclusively to the test set. No recording from that family was used for training, early stopping, learning-rate scheduling, or model selection. The recordings from each of the remaining thirteen ransomware families were independently divided into 80\% training and 20\% validation subsets using train\_\allowbreak{}test\_\allowbreak{}split \cite{ref46} with random\_\allowbreak{}state = 42. Performing the split separately within each family ensured that every in-distribution family with sufficient recordings was represented in both subsets.

Benign recordings were divided into 70\% training, 10\% validation, and 20\% test subsets using the same fixed random seed. The validation set was used exclusively for early stopping and learning-rate scheduling. The benign test subset and all recordings from the held-out malware family were not accessed during training or model selection.

The partitioning procedure was recording-disjoint but not binary-hash-disjoint or victim-set-disjoint within the thirteen in-distribution families. Consequently, recordings generated from the same malware binary or victim-file set could appear in both the training and validation subsets. The validation results therefore measure model-selection performance on new recordings of known families and potentially known variants; they should not be interpreted as evidence of within-family generalization to previously unseen binary variants. Cross-family generalization is evaluated only through the family-disjoint LOFO test set. A grouped sensitivity analysis using binary hash or acquisition provenance as the grouping key was not performed.

\subsection{Loss Function and Class Weighting}

The model is trained with BCEWithLogitsLoss (binary cross-entropy with logits). The positive class weight is set to min(max(1.0, ratio/\allowbreak{}2.0), 10.0), where ratio is the benign-to-malware process count ratio within the training partition of each fold. An upper bound of 10.0 was adopted after preliminary experiments showed that uncapped weights produced elevated false-positive rates. In conjunction with class weighting, per-process gradient accumulation normalizes the loss by the number of chunks, ensuring equal gradient contribution regardless of process length.

\subsection{Early Stopping and Learning Rate Schedule}

ReduceLROnPlateau monitors validation loss with patience = 1 and halves the learning rate (factor = 0.5) when no improvement is observed. Early stopping with patience = 3 terminates training \cite{ref47} when three consecutive epochs fail to improve validation loss. The best model checkpoint (lowest validation loss) is saved with the full state dictionary, optimizer state, validation loss, and configuration. Validation loss is monitored on the dedicated validation set, not the held-out test set, eliminating the data leakage present in the preliminary setup.

\section{Evaluation Methodology}

This part defines the evaluation metrics and protocols used in This thematic paper. It covers the primary metrics and the per-family evaluation procedure.

\subsection{Primary Metrics}

The primary operational metric is the process-level Missed Detection Rate (MDR), defined as MDR = FN /\allowbreak{} (TP + FN) $\times$ 100\%, where FN is the number of malware processes that are not detected and TP is the number of malware processes that are detected. A malware process is considered detected if its \textmu{}LSTM probability reaches or exceeds the fixed Stage 1 threshold of 0.30 in at least one chunk during its event stream. Otherwise, the process is counted as a false negative. This process-level definition assigns equal weight to each evaluated malware process, regardless of its total I/\allowbreak{}O volume.

A complementary run-level MDR is also reported where applicable. An execution run is considered detected if at least one malware process associated with that run reaches or exceeds the Stage 1 threshold. Consequently, a run containing multiple malware processes is counted as missed only when none of its eligible malware processes is detected. The process-level and run-level MDR values therefore represent different units of analysis and are reported separately.

Detection timeliness is measured using MBD3, defined as the 75th percentile of the cumulative WRITE bytes generated before the first detection among the detected malware processes. For each detected process, WRITE bytes are accumulated up to and including the first chunk whose \textmu{}LSTM probability reaches or exceeds 0.30. Bytes associated with READ, RENAME, and other non-WRITE events are excluded. Processes that never reach the threshold are represented by MDR and are not included in the MBD3 percentile calculation. A lower MBD3 value therefore indicates that detection occurred after less data had been written. MBD3 is reported in KiB, where 1 KiB equals 1,024 bytes.

The process-level MDR used in This thematic paper is not directly comparable with the MDR reported by CLEAR. CLEAR defines MDR as the proportion of missed ransomware traffic volume at a work point calibrated to a fixed false alarm rate, whereas This thematic paper measures the proportion of missed malware processes at a fixed probability threshold of 0.30. The two evaluations differ in both their units of analysis and their threshold-calibration procedures; therefore, their MDR and MBD3 values should be interpreted within their respective experimental settings rather than compared numerically.

\subsection{Per-Family Evaluation}

The retrained model is evaluated on all 14 families individually using the per-family evaluation pipeline. For each family, the evaluation loads up to 2,000 benign processes alongside all malware processes from that family, computes both streaming and batch metrics, and reports per-PID and per-run MDR, MBD3, AUC \cite{ref48}, and confusion matrix components.

\section{Real-Time Deployment Prototype}

This section presents a prototype implementation of the proposed detection system in a live setting. It is included to demonstrate practical feasibility and integration, rather than to define the core training and evaluation methodology of the thesis.

\subsection{System Design}

The real-time deployment prototype implements the complete detection pipeline as a single Python process with four cooperating threads: an event-ingestion thread (eBPF ring-buffer poll loop), an inference worker (per-PID tokenization and inference), a dashboard service (Flask/\allowbreak{}WebSocket state publisher), and a broadcast worker (1 Hz client push).

Each monitored process is tracked through a per-PID state object (PIDState) containing the token buffer, \textmu{}LSTM hidden state (h, c), I/\allowbreak{}O counters, and detection status (Figure 3.6). Process lifecycle tracking is implemented via a sched\_\allowbreak{}process\_\allowbreak{}exit tracepoint (primary) with a 30-second idle timeout as a fallback. The full thread role table, message queues, and shutdown protocol are given in Appendix E.1--E.2.

\begin{figure}[htbp]
\centering
\includegraphics[width=\linewidth]{fig3-6.png}
\caption{Per-PID state object structure and lifecycle management.}\label{fig:3-6}
\end{figure}

\subsection{3-Layer Detection Pipeline}

The two-stage detection architecture is implemented through three decision layers in the deployment prototype. Stage 1 is the \textmu{}LSTM sequence classifier. For deployment, Layer 1 applies a model-probability gate of 0.85 to each 100-event chunk. Stage 2 comprises two refinement layers: Layer 2 requires the probability gate to be satisfied for three consecutive chunks, and Layer 3 applies the heuristic score adjustments to the latest model probability. An alert is emitted only when the adjusted score remains at or above 0.85 after the consistency requirement has been met. The fixed threshold of 0.30 is used for the Stage 1 research metrics reported in the LOFO evaluation and is not the deployment alert threshold. Figure 3.7 illustrates the three deployment decision layers, and the rule weights and dashboard behavior are documented in Appendix E.3--E.4.

Based on the pipeline outcome, each monitored process is classified into one of four dashboard categories: ALERT (red --- quarantined), SUSPICIOUS (orange --- Layers 1 and 2 passed but Layer 3 inconclusive), SUPPRESSED (gray --- heuristic overrode the model), and NORMAL (green). Figure 3.7 illustrates the pipeline. Per-layer thresholds, rule score weights, and dashboard API details are documented in Appendix E.3--E.4.

\begin{figure}[htbp]
\centering
\includegraphics[width=\linewidth]{fig3-7.png}
\caption{The 3-layer real-time detection pipeline.}\label{fig:3-7}
\end{figure}

\chapter{Experimental Results}

\section{Chapter Overview}

This chapter evaluates the detection pipeline introduced in Chapter 3. It assesses cross-family generalization through leave-one-family-out (LOFO) evaluation across the fourteen ransomware families, analyzes the residual failure modes, and reports runtime and early-detection measurements from a separately trained deployment model that is not used to support generalization claims. All research metrics use the \textmu{}LSTM decision threshold of 0.30; the 0.85 alert threshold applies only to the deployment configuration (Sections 4.7 and 4.8).

\section{Evaluation Setup}

This part summarizes the experimental environment, the dataset partitions used for LOFO and prototype-deployment training, the fixed hyperparameter set adopted for the reported experiments, and the runs excluded from all reported metrics. The methodology underlying each item is described in Chapter 3; only the instantiated values used to produce the results reported in 4.3--4.5 and 4.7 are presented here. The hardware environments used for the reported experiments are summarized in Table 4.1

\begin{table}[htbp]
\centering
\small
\caption{Experimental hardware environments used for the reported results.}\label{tab:4-1}
\begin{tabularx}{\textwidth}{lXX}
\hline
\textbf{Environment} & \textbf{GPU configuration} & \textbf{Notes} \\
\hline
Kaggle & Tesla T4 $\times$2 (15 GB VRAM each) & 30 GB system RAM; 12-hour session limit \\
Google Colab Pro+ & Variable high-memory GPU environment (G4 \& L4) & Session-specific hardware allocation \\
\hline
\end{tabularx}
\end{table}

In all experiments, the filtered dataset was used, namely recording\_\allowbreak{}filtered.parquet for malware runs and the original Parquet files for benign recordings, together with precomputed WaR and RaR byte-overlap columns.

The fixed hyperparameter configuration used in all reported experiments is summarized in Table 4.2. These values were selected during preliminary development and frozen before any LOFO fold was evaluated. Because that preliminary development centered on the BlackBasta family, it is noted explicitly that the standard LOFO protocol applies to its fold as to every other: the held-out family's recordings appear only in the test partition, and early stopping uses validation loss computed over the remaining thirteen training families.

\begin{table}[htbp]
\centering
\small
\caption{Fixed hyperparameter set, frozen prior to all LOFO evaluations.}\label{tab:4-2}
\begin{tabularx}{\textwidth}{llX}
\hline
\textbf{Hyperparameter} & \textbf{Value} & \textbf{Purpose} \\
\hline
\multicolumn{3}{l}{\textbf{Model architecture}} \\
VOCAB & 1,024 & Token vocabulary size \\
EMBED & 128 & Embedding dimension \\
HIDDEN & 64 & LSTM hidden state size \\
LAYERS & 1 & Stacked LSTM layers \\
DROPOUT & 0.3 & Dropout before classifier head \\
\multicolumn{3}{l}{\textbf{Tokenization}} \\
CHUNK & 100 & Events per chunk (two 10-bit tokens per event) \\
\multicolumn{3}{l}{\textbf{Training}} \\
EPOCHS & 10 & Maximum training epochs \\
PATIENCE & 3 & Early-stop patience (epochs) \\
LR & 1$\times$10\textsuperscript{$-$}\textsuperscript{3} & Adam learning rate \\
GRAD\_\allowbreak{}CLIP & 1.0 & Gradient L2-norm clip threshold \\
MAX\_\allowbreak{}BENIGN\_\allowbreak{}PROCS & 5,000 & Per-fold benign process cap \\
\multicolumn{3}{l}{\textbf{Inference}} \\
THRESHOLD & 0.30 & \textmu{}LSTM decision threshold \\
\multicolumn{3}{l}{\textbf{Reproducibility}} \\
SEED & 42 & Random seed (numpy/\allowbreak{}torch/\allowbreak{}python) \\
\hline
\end{tabularx}
\end{table}

With the execution environment and training configuration fixed, Table 4.3 reports the binary diversity and recording-mode coverage of each evaluated ransomware family. The total run counts and captured event volumes are reported in Table 3.3 and are not repeated here.

\begin{table}[htbp]
\centering
\caption{Binary and recording-mode coverage by ransomware family.}\label{tab:4-3}
\begin{adjustbox}{max width=\textwidth}
\begin{tabular}{lrrr}
\hline
\textbf{Family} & \textbf{SHA-256 Binaries} & \textbf{Ransomware-Only Runs} & \textbf{Mixed-Mode Runs} \\
\hline
Akira & 11 & 72 & 65 \\
Bert & 2 & 10 & 10 \\
BlackBasta & 1 & 27 & 12 \\
Guunra & 2 & 10 & 19 \\
HelloKitty & 1 & 24 & 14 \\
IceFire & 1 & 15 & 20 \\
LockBit & 1 & 21 & 18 \\
Monti & 1 & 10 & 12 \\
Nitrogen & 1 & 20 & 15 \\
RAWorld & 1 & 5 & 5 \\
RansomEXX & 2 & 19 & 19 \\
Royal & 2 & 10 & 10 \\
Sodinokibi & 1 & 12 & 10 \\
SougoLock & 1 & 5 & 5 \\
Total & 28 & 260 & 234 \\
\hline
\end{tabular}
\end{adjustbox}
\end{table}

Note. Ransomware-only runs contain malware activity without the scripted benign workload, whereas mixed-mode runs contain malware activity executed concurrently with the scripted benign workload. The two recording modes comprise 494 malware-session runs in total.

Two low-signal recordings are excluded from all reported metrics: RUN\_\allowbreak{}IceFire\_\allowbreak{}Set1\_\allowbreak{}182048 and RUN\_\allowbreak{}MIXED\_\allowbreak{}Mirai\_\allowbreak{}Set1\_\allowbreak{}102810. Appendix D.3 documents these exclusions and the associated audit decisions.

The per-family run count ranges from 10 (RAWorld, SougoLock) to 137 (Akira) --- a 13.7$\times$ spread that reflects three concurrent factors. First, every distinct Linux-targeting binary variant available per family was collected: Akira had 11 distinct hash variants whereas SougoLock and RAWorld each had only one. Second, the malware orchestrator targeted 15 detonation rounds per family, so multi-hash families naturally accumulated more successful runs across their variants. Third, a small fraction of planned detonations did not produce a usable recording because the binary crashed or exceeded the 30-minute wall-clock budget; only successful executions are counted. The run counts were not artificially balanced after data collection: truncating to the minimum would discard cross-variant diversity in well-represented families --- notably the 11 Akira variants, including the false-negative-prone hash analyzed in Section 4.5.1 --- while oversampling would introduce duplicate trajectories that the per-process LOFO evaluation already treats as independent test points.

\section{Leave-One-Family-Out (LOFO) Results}

This part presents the main cross-family generalization results under the LOFO protocol. Each of the fourteen families is held out in turn, with training performed on the remaining thirteen and evaluation conducted on the held-out family using streaming maximum-probability detection. The reported metrics (MDR, AUC, recall, precision, and MBD3) follow the per-process, fixed-threshold definitions introduced in Section 3.7 and are therefore not directly comparable to CLEAR's volume-weighted, false-alarm-rate-calibrated results. This qualification is revisited in Section 4.9 and Section 5.2.

A structural feature of the dataset bears on how the LOFO folds should be read, and the two fold types are therefore reported separately rather than pooling them. Five families --- Akira, Bert, Guunra, RansomEXX, and Royal --- contribute between two and eleven distinct SHA-256 binaries each (nineteen hashes in total), so their held-out folds test generalization across binary variants the model has not seen, that is, cross-family generalization in the intended sense. The remaining nine families contribute exactly one binary each; for those folds the held-out set is a single binary, so the leave-one-family-out protocol reduces to a leave-one-binary-out experiment and cross-family generalization cannot be separated from binary-fingerprint generalization. Of the ten folds reaching MDR = 0\%, four fall in the multi-binary group (Bert, Guunra, RansomEXX, Royal) and six in the single-binary group; the pooled 10.49\% missed-detection rate should be read against this split. Acquiring additional variants per family and reporting a paired leave-one-variant-out evaluation alongside LOFO is identified as future work in Section 5.4.1.

\begin{table}[htbp]
\centering
\caption{Leave-One-Family-Out detection performance (14 families).}\label{tab:4-4}
\begin{adjustbox}{max width=\textwidth}
\footnotesize
\begin{tabular}{lllllrrrl}
\hline
\textbf{Family} & \textbf{MDR (\%)} & \textbf{AUC} & \textbf{Recall} & \textbf{Precision} & \textbf{TP} & \textbf{FN} & \textbf{FP} & \textbf{MBD3 (KiB)} \\
\hline
\multicolumn{9}{l}{\textbf{Group 1: Encryption-dominant}} \\
Bert & 0.00 & 1.0000 & 1.0000 & 0.5128 & 20 & 0 & 19 & 1956.8 \\
BlackBasta & 0.00 & 1.0000 & 1.0000 & 0.7500 & 39 & 0 & 13 & 33.1 \\
HelloKitty & 34.48 & 0.7190 & 0.6552 & 0.8444 & 38 & 20 & 7 & 10180.7 \\
IceFire & 0.00 & 1.0000 & 1.0000 & 0.6863 & 35 & 0 & 16 & 129.2 \\
Nitrogen & 0.00 & 1.0000 & 1.0000 & 0.7292 & 35 & 0 & 13 & 9407.7 \\
Sodinokibi & 4.35 & 0.9914 & 0.9565 & 0.5641 & 22 & 1 & 17 & 4182.5 \\
SougoLock & 100.00 & 0.9450 & 0.0000 & 0.0000 & 0 & 10 & 4 & N/\allowbreak{}A \\
\multicolumn{9}{l}{\textbf{Group 2: Rename-dominant}} \\
Guunra* & 0.00 & 1.0000 & 1.0000 & 0.7838 & 29 & 0 & 8 & 0.0 \\
\multicolumn{9}{l}{\textbf{Group 3: Bimodal}} \\
Akira & 16.79 & 0.9975 & 0.8321 & 0.8085 & 114 & 23 & 27 & 13.2 \\
LockBit & 0.00 & 0.9999 & 1.0000 & 0.6290 & 39 & 0 & 23 & 103.7 \\
Monti & 0.00 & 1.0000 & 1.0000 & 0.8148 & 22 & 0 & 5 & 16421.5 \\
RAWorld & 0.00 & 1.0000 & 1.0000 & 0.3571 & 10 & 0 & 18 & 447949.4 \\
RansomEXX & 0.00 & 1.0000 & 1.0000 & 0.7600 & 38 & 0 & 12 & 16930.9 \\
Royal & 0.00 & 1.0000 & 1.0000 & 0.6452 & 20 & 0 & 11 & 5230.8 \\
\hline
\end{tabular}
\end{adjustbox}
\end{table}

Note. Families are ordered by behavioral cluster rather than alphabetically. Group 1 = encryption-dominant, Group 2 = rename-dominant, and Group 3 = bimodal. MDR and MBD3 follow the per-process, fixed-threshold definitions introduced in Chapter 3. * Guunra detection signal is limited to RENAME events; WRITE operations performed via pwritev/\allowbreak{}preadv/\allowbreak{}io\_\allowbreak{}uring are not captured by the current sensor (see Section 4.5.2). The reported MDR therefore measures the model's ability to detect on the RENAME stream alone.

Across the fourteen held-out-family folds, twelve families achieved an AUC of 0.99 or higher. The two exceptions were HelloKitty (AUC = 0.7190) and SougoLock (AUC = 0.9450), both of which exhibited systematic failure modes analyzed in Section 4.5. AUC measures ranking performance across all possible thresholds and should not be interpreted as detection performance at the fixed Stage 1 evaluation threshold of 0.30. This distinction is particularly important for SougoLock: although its AUC was 0.9450, all ten malware processes remained below the Stage 1 evaluation threshold, resulting in Recall = 0 and MDR = 100\%.

Across the fourteen held-out-family folds, the pooled per-process MDR was 10.49\% among 515 malware processes. A 2,000-resample percentile bootstrap produced a 95\% interval of [7.96\%, 13.01\%]. This bootstrap resampled individual malware processes from the pooled folds and calculated the proportion with a maximum probability below the fixed Stage 1 evaluation threshold of 0.30. The mean per-fold AUC was 0.9752, with a 95\% bootstrap interval of [0.9676, 0.9825]. For this interval, positive and negative processes were resampled separately within each fold, the fold AUC values were recalculated, and their mean was recorded for each bootstrap sample.

These intervals quantify uncertainty under process-level resampling of the existing predictions. They treat processes as independent observations and do not account for dependence among processes originating from the same recording, execution run, or binary hash. They also do not measure variability arising from alternative data partitions, training seeds, or model retraining. The intervals should therefore be interpreted as conditional descriptive intervals for the observed predictions rather than as evidence that the reported performance is stable across independently sampled runs or ransomware variants. A run- or hash-clustered bootstrap was not performed and remains necessary for stronger population-level inference. Figure 4.1 presents the per-family MDR values across the fourteen LOFO folds.

\begin{figure}[htbp]
\centering
\includegraphics[width=\linewidth]{fig4-1.png}
\caption{Per-family MDR across LOFO folds.}\label{fig:4-1}
\end{figure}

\section{Per-Family Analysis}

This section characterizes the VFS-level behavioral patterns observed across the fourteen ransomware families. The grouping is based on the recorded event distributions and is independent of whether a family was detected successfully at the fixed Stage 1 evaluation threshold. Systematic detection failures and their causes are examined separately in Section 4.5.

\subsection{Behavioral Groups}

The fourteen families are assigned to three behavioral groups based on the per-family I/\allowbreak{}O profiles introduced in Section 3.3.1: encryption-dominant, rename-dominant, and bimodal. Table 4.5 reports the assigned group, the observed WaR and RaR frequencies, and the median write volume per run. These measurements describe the VFS events visible to the instrumented hooks and provide context for the auxiliary-feature analysis reported in Appendix F.

The encryption-dominant group contains families whose captured traces are primarily characterized by sustained or repeated WRITE activity. The observed patterns range from small fixed-chunk in-place writes, as exhibited by BlackBasta, to larger bulk and medium-chunk writes observed in families such as Bert, IceFire, and Nitrogen. Membership in this group describes the dominant captured behavior and does not imply uniformly successful detection.

The rename-dominant group contains Guunra, for which the instrumented hooks captured a predominantly RENAME-based trace and no WRITE events. This classification describes the events visible to the current sensor and does not imply that no underlying file-content modification occurred. Detection of this family appears to benefit from the categorical RENAME signal encoded in the token representation.

The bimodal group contains families whose recordings exhibit two distinguishable behavioral modes, either within individual runs or across runs from the same family. These modes may include alternating write patterns, rename activity, and low-activity intervals. Akira, LockBit, Monti, RAWorld, RansomEXX, and Royal are assigned to this group. The group label refers to the family-level distribution of captured behavior and should not be interpreted as a mixed-workload condition.

\begin{table}[htbp]
\centering
\caption{Behavioral groups with per-family WaR/\allowbreak{}RaR frequencies and median write volume.}\label{tab:4-5}
\begin{adjustbox}{max width=\textwidth}
\small
\begin{tabular}{llrlll}
\hline
\textbf{Family} & \textbf{Group} & \textbf{Runs} & \textbf{WaR (\%)} & \textbf{RaR (\%)} & \textbf{Bytes/\allowbreak{}run (median)} \\
\hline
Bert & Encryption-dominant & 20 & 45.04 & 69.99 & 73.05 GB \\
BlackBasta & Encryption-dominant & 39 & 97.62 & 96.91 & 150.20 MB \\
HelloKitty & Encryption-dominant & 38 & 59.46 & 56.81 & 23.60 GB \\
IceFire & Encryption-dominant & 35 & 19.95 & 0.02 & 63.98 MB \\
Nitrogen & Encryption-dominant & 35 & 56.30 & 26.14 & 51.95 GB \\
Sodinokibi & Encryption-dominant & 22 & 45.30 & 66.96 & 48.77 GB \\
SougoLock & Encryption-dominant & 10 & 0.05 & 49.73 & 12.13 GB \\
Guunra & Rename-dominant & 29 & 0.00 & 0.00 & 0 \\
Akira & Bimodal & 137 & 31.16 & 0.09 & 204.97 KB \\
LockBit & Bimodal & 39 & 49.60 & 74.14 & 159.20 GB \\
Monti & Bimodal & 22 & 54.14 & 0.15 & 41.49 GB \\
RAWorld & Bimodal & 10 & 33.32 & 0.00 & 62.54 GB \\
RansomEXX & Bimodal & 38 & 36.18 & 1.50 & 28.86 GB \\
Royal & Bimodal & 20 & 28.80 & 0.00 & 77.13 GB \\
\hline
\end{tabular}
\end{adjustbox}
\end{table}

Note. WaR and RaR are reported as percentages of the corresponding VFS events. *Guunra produced no WRITE events through the instrumented hooks; the zero value does not indicate that no file modification occurred.

\subsection{Early-Detection Traces}

This subsection illustrates the progression of Stage 1 detection probabilities using two representative process traces. In each figure, the horizontal axis represents consecutive chunks of one hundred VFS events, the primary vertical axis represents the \textmu{}LSTM probability assigned to each chunk, and the horizontal reference line marks the fixed Stage 1 evaluation threshold of 0.30. The secondary vertical axis reports cumulative WRITE bytes. A process is considered detected at the first chunk whose probability reaches or exceeds the threshold. The cumulative WRITE bytes observed up to that point provide the process-level value used in the MBD3 calculation.

Figure 4.2 presents an immediate-detection trace from BlackBasta. The probability reaches the threshold in the first chunk, after approximately 33 KiB of cumulative writes. Figure 4.3 presents a burst-structured trace from Nitrogen. Detection occurs in the first chunk, after approximately 639 KiB of cumulative writes. Two periods of elevated activity are separated by a low-activity interval during which both the write volume and model probability decrease. This temporal pattern does not change Nitrogen's family-level assignment to the encryption-dominant group.

\begin{figure}[htbp]
\centering
\includegraphics[width=\linewidth]{fig4-2.png}
\caption{Encryption-dominant example (BlackBasta).}\label{fig:4-2}
\end{figure}

\begin{figure}[htbp]
\centering
\includegraphics[width=\linewidth]{fig4-3.png}
\caption{Burst-structured detection trace from Nitrogen.}\label{fig:4-3}
\end{figure}

\section{Error Analysis}

This section examines the systematic failure modes identified during evaluation, evidenced in Table 4.6 (Section 4.5.1), Table 4.7 (Section 4.5.2), and Table F.5 (Section 4.5.3). The three subsections are independent in root cause and appear in decreasing order of impact on reported metrics.

\subsection{Log-Spam Dilution (Akira 0ee1d284ed663073)}

Among the 137 Akira test processes in the leave-one-family-out fold, 23 processes are missing at the 0.30 decision threshold. A per-process profile of these false negatives, computed against the 112 correctly detected Akira processes, reveals a clean two-population separation. The missed processes are characterized by extreme path concentration: the median fraction of events directed at a single path is 97.34\%, with 92\% of the missed processes meeting a ``log-only writer'' criterion (more than 80\% of events on the top path). The event mix of the missed group was measured at 99.20\% writes and 0.27\% renames, versus 0.03\% writes and 99.96\% renames in the detected group. The detected processes also exhibit broader path coverage, with a median of 95,993 distinct paths touched per process versus 7,102 for the missed group.

This pattern indicates a sequence-level dilution effect rather than a tokenization failure. The Akira hash 0ee1d284ed663073 (and several others in the family) writes an abnormally large volume of operational log entries while encrypting. Each log write is correctly tokenized as a WRITE-aux combination, but the resulting token sequence is dominated by these repeated log tokens, which the \textmu{}LSTM integrates into its hidden state and which mask the lower-frequency encryption tokens within the same per-process sequence. The suppressed peak probability for this hash, relative to the mean curve for other Akira hashes, indicates that the signal is present in the token stream but is sequence-statistically dominated by log-write tokens; the per-chunk heuristic gate (Section 3.8.2) cannot recover the signal because the dilution occurs within rather than across chunks.

The broad-cap retrain was correlational: capping every path at 100 events recovered most Akira misses, but the intervention simultaneously changed total token volume, sequence length, class balance, and log-path dominance --- leaving the causal mechanism unidentified. Setup A isolates the log-spam factor. Within each PID, the path whose event share exceeds 80\% is identified, and tokens are capped only on that path, leaving every other path uncapped. All other training-time variables --- seed, fold split, hyperparameters, benign sub-pool, stateful state carry, WaR/\allowbreak{}RaR semantics --- are held identical to baseline. The result on the Akira fold is MDR~=~0.00\% with all 137 Akira processes detected (Table 4.6); the broad-cap retrain achieved only MDR~=~5.84\%. Because Setup A varies only the log-path-only axis and yet matches or exceeds the broad cap, total volume, sequence length, and class balance are ruled out as the active mechanism. These results identify log-spam dilution as the dominant factor for the Akira false-negative cohort; the isolation rests on a single-seed controlled run on one fold and is not multi-seed validated.

\begin{table}[htbp]
\centering
\small
\caption{Akira LOFO fold: three-way controlled ablation isolating the log-spam factor (Setup A).}\label{tab:4-6}
\begin{tabularx}{\textwidth}{Xlllllrrrl}
\hline
\textbf{Configuration} & \textbf{MDR (\%)} & \textbf{AUC} & \textbf{Precision} & \textbf{Recall} & \textbf{F1} & \textbf{TP} & \textbf{FN} & \textbf{FP} & \textbf{MBD3 (KiB)} \\
\hline
Baseline (no cap) & 16.79 & 0.9975 & 0.8085 & 0.8321 & 0.8201 & 114 & 23 & 27 & 13.2 \\
Broad cap~=~100 (all paths) & 5.84 & 0.9998 & 0.9699 & 0.9416 & 0.9555 & 129 & 8 & 4 & --- \\
Setup A: log-path-only cap~=~100 & 0.00 & 0.9999 & 0.8107 & 1.0000 & 0.8954 & 137 & 0 & 32 & 70.6 \\
\hline
\end{tabularx}
\end{table}

Note. All counts on held-out Akira test processes (n = 137); only the capping strategy varies (seed = 42, same fold, frozen hyperparameters). MBD3 = P75 cumulative WRITE bytes before first detection; --- = not detected

\subsection{Sensor Blind Spots (Guunra and Akira pwritev variants)}

The eBPF sensor (Section 3.2.2) hooks vfs\_\allowbreak{}read, vfs\_\allowbreak{}write, and vfs\_\allowbreak{}rename. The real-encryptor variants of Guunra and Akira bypass the read/\allowbreak{}write hooks by performing their file I/\allowbreak{}O through the vectored readv/\allowbreak{}writev/\allowbreak{}pwritev system-call family, which the kernel dispatches via vfs\_\allowbreak{}iter\_\allowbreak{}read/\allowbreak{}vfs\_\allowbreak{}iter\_\allowbreak{}write rather than the scalar vfs\_\allowbreak{}read/\allowbreak{}vfs\_\allowbreak{}write functions the kprobes attach to --- a documented evasion vector for I/\allowbreak{}O-monitoring-based detection \cite{ref24}. Each mechanism was confirmed by running the binaries under strace in the sandbox VM and verifying with an md5 checksum that the file content was genuinely rewritten, not merely renamed. Akira encrypts in place, delivering its ciphertext through pwritev() in chunks at increasing offsets; Guunra reads each file whole through readv() and writes the ciphertext back through writev(), followed by a trailing key blob. In both cases the read of the plaintext and the write of the ciphertext are invisible to the sensor, and the run therefore appears as a RENAME-only trail --- but the subsequent rename() to the ransom extension is captured by the vfs\_\allowbreak{}rename hook, which is why detection still succeeds (Guunra LOFO AUC = 1.0, via the rename-only token pattern). This is a feature-richness limitation, not a detection-capability one: the WRITE-opcode signal is lost, but the rename signature is not. Extending the sensor to hook the vectored-I/\allowbreak{}O path and io\_\allowbreak{}uring submission queues is listed as Future Work in Chapter 5.

The three families with systematic failure modes exhibit distinct event-mix signatures arising from three separate mechanisms. Guunra and the eight pwritev-based Akira variants listed in Table 4.7 are both genuine vectored-I/\allowbreak{}O blind spots: their captured runs are dominated by RENAME with almost no visible WRITE, because the encryption writes route through bypassed paths and never reach the sensor. The leave-one-family-out fold still detects both via the rename-only token pattern alone. Akira additionally exhibits a second, unrelated failure mode at the process level: among the false-negative processes analyzed in Section 4.5.1, a single high-frequency path accounts for over 97\% of events (99\% writes, 0.27\% renames)---a log-spam dilution effect, not a sensor blind spot, since these writes are fully visible to the sensor but are diluted within the token sequence. The SougoLock profile reflects a separate, third limitation: the offset-resolution limit of Section 4.5.3, where the coarse 4-kilobyte offset granularity collapses distinct write locations into the same token. Table 4.7 reports the per-hash opcode breakdown for the Akira and Guunra vectored-I/\allowbreak{}O blind-spot variants.

\begin{table}[htbp]
\centering
\caption{Per-hash opcode breakdown for the Akira and Guunra runs affected by vectored-I/\allowbreak{}O sensor blind spots.}\label{tab:4-7}
\begin{adjustbox}{max width=\textwidth}
\small
\begin{tabular}{llrrlll}
\hline
\textbf{Family} & \textbf{Hash} & \textbf{Runs} & \textbf{Events} & \textbf{W \%} & \textbf{R \%} & \textbf{RN \%} \\
\hline
Akira & 12ac468c & 10 & 823,904 & 0.1\% & 0.0\% & 99.9\% \\
 & 1bf45e64 & 24 & 1,765,172 & 0.0\% & 0.0\% & 100.0\% \\
 & 1d3b5c65 & 9 & 795,212 & 0.1\% & 0.0\% & 99.9\% \\
 & 718d6f18 & 10 & 1,057,386 & 0.0\% & 0.0\% & 100.0\% \\
 & 7454191b & 10 & 1,057,386 & 0.0\% & 0.0\% & 100.0\% \\
 & aacc05b2 & 10 & 1,057,386 & 0.0\% & 0.0\% & 100.0\% \\
 & acfe720d & 10 & 1,058,246 & 0.1\% & 0.0\% & 99.9\% \\
 & bcae978c & 10 & 1,058,246 & 0.1\% & 0.0\% & 99.9\% \\
Guunra & 5677dfad & 14 & 3,791,602 & 0.0\% & 0.0\% & 100.0\% \\
 & 75cb7eb7 & 15 & 3,730,451 & 0.0\% & 0.0\% & 100.0\% \\
Total & --- & 122 & 16,194,991 & 0.0\% & 0.0\% & 100.0\% \\
\hline
\end{tabular}
\end{adjustbox}
\end{table}

Note. Two additional Akira variants (3b4137e8, 5114c107) rename files to the .akira extension but do not encrypt content on Linux (md5-confirmed unchanged); they are ESXi-oriented builds whose encryption path is gated to that platform. They are excluded from this table because they do not exhibit the vectored-I/\allowbreak{}O write pattern analyzed here.

\subsection{Root Causes of Token-Representation Failures}

Two token-representation limitations were observed in the experimental evaluation: coarse offset resolution in the SougoLock fold and limited write-size granularity in a subset of Akira runs. The controlled ablation supports an offset-resolution explanation only within the evaluated SougoLock fold; the Akira write-size observation is descriptive and was not isolated from classifier effects. All ten SougoLock test processes were missed (MDR = 100\%, recall = 0), with predicted malware probabilities clustered near 10\textsuperscript{$-$}\textsuperscript{6}, far below the 0.30 detection threshold.

To isolate the cause, a controlled ablation retrained the same fold under five tokenizer configurations, holding all other settings identical (Table F.5). Two candidate causes were ruled out: recording every read interval per path instead of only the latest one reduced the miss rate only from 100\% to 90\%, and removing the delta-time feature entirely left the miss rate unchanged at 100\%. Among the five configurations evaluated, only widening the offset field from six to twelve bits closed the gap completely. This increased the effective offset resolution from 4 KiB to 256 bytes, reducing the miss rate to 0.0\% and raising the fold AUC from 0.9450 to 0.9990.

Within the evaluated SougoLock fold, the results identify insufficient tokenizer offset resolution as the principal cause of the observed detection failure. Expanding the offset allocation changed the outcome from complete failure to full detection, whereas disabling the model's cross-chunk memory left the baseline failure unchanged. This controlled ablation provides diagnostic evidence that representation resolution was the limiting factor in this specific fold. However, because the expanded representation was evaluated only on the SougoLock fold and increased the false-positive count from 4 to 23, the result does not establish that representation is the dominant determinant of detection performance across all fourteen families.

The twelve-bit configuration is presented as a diagnostic finding rather than a production change: it increased false positives on this fold from four to twenty-three and has not been retrained across all fourteen families. The production tokenizer therefore retains the six-bit offset, with a budget-aware twelve-bit retraining left as future work.

While the SougoLock case demonstrates a coarse offset-resolution failure, a separate limitation was observed in Akira: a subset of runs generated single-byte writes that the current four-bit logarithmic size field maps to the same low-end size bin as other small writes. This issue concerns write-size quantization rather than offset width; the twelve-bit offset configuration does not address it. Potential remedies include allocating additional bits to the size field or adding a supplementary write-size distribution feature.

\section{Heuristic False-Positive Suppression}

This section reports the heuristic rule layer, a deterministic post-filter evaluated for false-positive suppression downstream of the \textmu{}LSTM. Appendix F reports diagnostic ablations of output frequency, decision aggregation, statefulness, and the Token-2 auxiliary-bit group. The contribution of the auxiliary-bit group was family dependent and substantial in the BlackBasta fold.

The heuristic rule layer is evaluated here as an in-sample diagnostic on the observed benign false-positive pool. In this analysis, the false-positive count decreases from 193 to 4 across the fourteen LOFO folds (97.9\%), and mean per-family F1 increases from 0.7274 to 0.9016. These F1 values retain the Stage 1 TP and FN counts; therefore, Table 4.8 quantifies false-positive suppression but does not independently establish true-positive retention. The diagnostic uses the Stage 1 research threshold of 0.30 to define candidate processes and should not be interpreted as the deployment decision order. In the deployment prototype, an unadjusted model probability of at least 0.85 is first required for three consecutive chunks, after which the heuristic checker adjusts the latest probability and an alert is emitted only if the adjusted score remains at least 0.85 (Section 3.8.2 and Appendix E.3). Table 4.8 summarizes the diagnostic comparison.

\begin{table}[htbp]
\centering
\caption{Heuristic ablation across the fourteen LOFO folds.}\label{tab:4-8}
\begin{adjustbox}{max width=\textwidth}
\footnotesize
\begin{tabular}{lrrlrrll}
\hline
\textbf{Family} & \textbf{TP} & \textbf{FN} & \textbf{MDR (\%)} & \textbf{FP (S1)} & \textbf{FP (S2 opt)} & \textbf{F1 (S1)} & \textbf{F1 (S2 opt)} \\
\hline
Akira & 114 & 23 & 16.79 & 27 & 0 & 0.820 & 0.908 \\
Bert & 20 & 0 & 0.00 & 19 & 0 & 0.678 & 1.000 \\
BlackBasta & 39 & 0 & 0.00 & 13 & 0 & 0.857 & 1.000 \\
Guunra & 29 & 0 & 0.00 & 8 & 0 & 0.879 & 1.000 \\
HelloKitty & 38 & 20 & 34.48 & 7 & 0 & 0.738 & 0.792 \\
IceFire & 35 & 0 & 0.00 & 16 & 2 & 0.814 & 0.972 \\
LockBit & 39 & 0 & 0.00 & 23 & 0 & 0.772 & 1.000 \\
Monti & 22 & 0 & 0.00 & 5 & 0 & 0.898 & 1.000 \\
Nitrogen & 35 & 0 & 0.00 & 13 & 2 & 0.843 & 0.972 \\
RAWorld & 10 & 0 & 0.00 & 18 & 0 & 0.526 & 1.000 \\
RansomEXX & 38 & 0 & 0.00 & 12 & 0 & 0.864 & 1.000 \\
Royal & 20 & 0 & 0.00 & 11 & 0 & 0.784 & 1.000 \\
Sodinokibi & 22 & 1 & 4.35 & 17 & 0 & 0.710 & 0.978 \\
SougoLock & 0 & 10 & 100.00 & 4 & 0 & 0.000 & 0.000 \\
Aggregate & --- & --- & --- & 193 & 4 & 0.7274 & 0.9016 \\
\hline
\end{tabular}
\end{adjustbox}
\end{table}

\section{Production Model Configuration and Inference Microbenchmark}

For the deployment measurements that follow, a separately trained prototype model is used. Unlike the LOFO folds, this model is trained on an 85/\allowbreak{}10/\allowbreak{}5 random split of all fourteen families, with the explicit goal of reporting model-inference latency, throughput, memory footprint, and early-detection behavior under a representative production configuration. Because this model has seen every family during training, it does not support cross-family generalization claims; those claims remain exclusive to Section 4.3. The measurements that follow characterize deployment feasibility --- latency, throughput, and memory footprint under a representative production configuration --- rather than the thesis's core detection-quality contributions, which are evaluated separately in Sections 4.3 and 4.4. Table 4.9 summarizes the training configuration and in-training validation results for this deployment model.

\begin{table}[htbp]
\centering
\small
\caption{Final deployment model: training configuration and in-training validation results.}\label{tab:4-9}
\begin{tabularx}{\textwidth}{XX}
\hline
\textbf{Parameter} & \textbf{Value} \\
\hline
Train set (files) & 571 ($\approx$ 85 \%) \\
Validation set (files) & 72 ($\approx$ 10 \%) \\
Test set (files) & 39 ($\approx$ 5 \%) \\
Total (files) & 682 \\
Stratification & Per-family malware + random benign; seed = 42 \\
Architecture and hyperparameters & As in Table 4.2 \\
Parameters & 182,849 \\
Best validation loss & 0.0086 \\
Val loss at epoch 10 & 0.0194 (overfit; best ckpt at epoch 7) \\
Train loss at epoch 10 & 0.0008 \\
Pos-weight (auto) & 4.34 \\
Epochs run & 10 /\allowbreak{} 10 (full budget) \\
Best epoch (deployed ckpt) & 7 (val\_\allowbreak{}loss = 0.0086) \\
Training time (Colab T4) & $\approx$ 11 hours ($\approx$ 66 min/\allowbreak{}epoch $\times$ 10) \\
Seed & 42 \\
Total processes (test set) & 71,073 \\
MDR (\%) & 0.00 \\
Recall & 1.0000 \\
AUC & 0.9998 \\
\textmu{}LSTM FP (test set) & 788 (1.11 \%) \\
Checkpoint file & clear\_\allowbreak{}ulstm\_\allowbreak{}merged\_\allowbreak{}final.pth \\
Checkpoint size & 718.6 KiB (735,850 bytes) \\
\hline
\end{tabularx}
\end{table}

Training ran the full ten-epoch budget; validation loss reached its minimum of 0.0086 at epoch 7 before rising to 0.0194 at epoch 10, a pattern consistent with overfitting after the best checkpoint was saved. The best-checkpoint mechanism preserved the epoch-7 model as the production checkpoint.

\begin{table}[htbp]
\centering
\small
\caption{Isolated CPU inference microbenchmark on the Windows 11 development host.}\label{tab:4-10}
\begin{tabularx}{\textwidth}{llX}
\hline
\textbf{Group} & \textbf{Parameter} & \textbf{Value} \\
\hline
Hardware & Measurement host & Windows 11 development host, single-thread CPU inference \\
Hardware & Deployment target & Ubuntu 24.04.4, kernel 6.17 (VM guest: 14 vCPU, 8 GB RAM) \\
Latency (CPU) & Mean per chunk & 0.855 ms \\
Latency (CPU) & Median per chunk & 0.838 ms \\
Latency (CPU) & P95 per chunk & 1.143 ms \\
Latency (CPU) & P99 per chunk & 1.743 ms \\
Latency (CPU) & Min /\allowbreak{} max per chunk & 0.624 /\allowbreak{} 2.127 ms \\
Throughput (CPU) & Chunks /\allowbreak{} second & 1,169 \\
Throughput (CPU) & Events /\allowbreak{} second & 116,914 \\
Memory (1,021 PIDs) & Baseline RSS & 309.3 MB \\
Memory (1,021 PIDs) & Peak RSS & 435.9 MB \\
Memory (1,021 PIDs) & Growth (n = 1,021 procs) & 126.6 MB \\
Memory (1,021 PIDs) & Model state on disk & 0.7 MB \\
\hline
\end{tabularx}
\end{table}

A supplementary inference measurement was performed inside the Ubuntu 24.04 VirtualBox guest used for malware capture. The median inference latency was 0.43 ms in the Linux VM and 0.84 ms on the Windows host. However, the Linux VM exhibited a P99 latency of 327 ms, compared with 1.74 ms on the Windows host. The difference may reflect host contention, virtualization overhead, scheduling variability, or other uncontrolled system activity. Because the measurement did not isolate these factors, the Windows result is retained as the canonical model-inference microbenchmark, while the Linux VM result demonstrates that tail latency may differ substantially in the deployment environment.

The model contains 182,849 trainable parameters and occupies approximately 0.7 MB on disk. On the Windows development host, the isolated inference loop achieved a median latency of 0.84 ms per chunk, with P95 and P99 values of 1.14 ms and 1.74 ms, respectively, and a measured throughput of 116,914 events per second. These results indicate that the model inference operation itself is computationally lightweight on the tested CPU. They do not establish end-to-end latency because upstream event capture and downstream alert-processing delays were not included.

The observed resident-set size increased by 126.6 MB while 1,021 PIDs were processed. The theoretical LSTM hidden and cell tensors require approximately 512 bytes per active PID, corresponding to approximately 0.5 MB for 1,021 PIDs. The remaining resident-set growth cannot be attributed to the hidden state alone and may include Python objects, per-process event buffers, path data, interpreter overhead, and memory-allocator caching. Because component-level memory profiling was not performed, the 126.6 MB increase is reported as total observed process-memory growth rather than as the memory cost of the \textmu{}LSTM hidden state.

\section{Live-VM Validation}

\subsection{Motivation and Test Environment}

The experiments in this section evaluate whether the pipeline operates correctly in the Ubuntu deployment environment. The five tests have different evidential roles: T1, T2, and T4 exercise the live sensor with benign workloads; T3 and T5 replay previously recorded Parquet traces through the inference pipeline; and the mixed-workload experiment in Section 4.8.3 executes a BlackBasta binary using its held-out LOFO checkpoint. The tests were conducted in a VirtualBox virtual machine running Ubuntu 24.04.4 with kernel 6.17 and BCC 0.29.1. The deployment model contains 182,849 parameters, and the live sensor attaches kprobes to vfs\_\allowbreak{}read, vfs\_\allowbreak{}write, and vfs\_\allowbreak{}rename. The alert configuration applies a threshold of 0.85 sustained for three consecutive chunks, followed by heuristic verification. These tests assess pipeline integration under the evaluated environment but do not provide equivalent evidence of held-out-family detection.

\subsection{Tests and Observations}

The five deployment and integration tests are summarized in Table 4.11. The dashboard was restarted before each test, and the alert log was inspected after the test window. Tests T1, T2, and T4 generated live activity through the attached eBPF sensor, whereas T3 and T5 processed previously recorded Parquet traces through the inference pipeline.

\begin{table}[htbp]
\centering
\small
\caption{Summary of the five deployment and integration tests on Ubuntu 24.04.4 (kernel 6.17).}\label{tab:4-11}
\begin{tabularx}{\textwidth}{lXlrX}
\hline
\textbf{Test} & \textbf{Workload} & \textbf{Sensor capture} & \textbf{Alerts} & \textbf{Headline observation} \\
\hline
T1 & 30-second idle dashboard & 0 events & 0 & pipeline boots; probes attach \\
T2 & 30 git commits + 30 file copies & 2,838 events & 0 & benign with trusted-name workload \\
T3 & Akira parquet replay (PID 2538) & 19,937 /\allowbreak{} 200 chunks & 1 & bytes to first detection = 0.326 KiB \\
T4 & 6 untrusted-comm workloads (dd /\allowbreak{} openssl /\allowbreak{} sha256sum /\allowbreak{} shred /\allowbreak{} make /\allowbreak{} gcc) & 2,898 events & 0 & selection-bias check --- 0 alerts \\
T5 & SougoLock parquet replay (PID 2679) & 26,498 /\allowbreak{} 265 chunks & 1 & bytes to first detection $\approx$ 159.5 MiB \\
\hline
\end{tabularx}
\end{table}

Note. Sensor capture is the count of VFS events observed by the eBPF probes; alerts is the count of entries written to the alert log during the test window.

Test T1 is a sanity check: with no workload generated, the sensor captured zero events and the alert log remained empty for the full thirty-second window.

Test T2 exercises a benign workload of thirty git commits and thirty file copies in a dedicated directory. 2,838 captured events and zero alerts were measured. Because git is in the allowlist, T2 alone does not demonstrate generalization beyond the allowlist, a question addressed by T4.

Test T3 was an offline replay integration test using a previously recorded Akira trace. The trace produced an alert at chunk 0 after 0.326 KiB of cumulative writes. Because the deployment model had been trained using data from all fourteen families, including Akira, this result demonstrates that the inference and alerting components operated correctly in the Linux environment but does not provide evidence of held-out-family detection or generalization.

Test T4 exercised six benign workloads whose process names were absent from the sixty-eight-name trusted-process allowlist: dd, openssl, sha256sum, shred, make, and gcc. No alerts were recorded across the 2,898 captured events. This result shows that the complete pipeline did not produce false alerts for these six workloads under the tested configuration. However, the small workload set and the absence of repeated trials prevent the result from establishing generalization to unseen benign workloads or an end-to-end false-positive rate of zero.

Test T5 was an offline replay integration test using a previously recorded SougoLock trace. The replay produced an alert at chunk 0 after approximately 159.5 MiB had already been written. Because the deployment model had been trained using SougoLock data, the alert represents an in-sample replay result. It does not alter the held-out SougoLock result, which remains an MDR of 100\% at the Stage 1 threshold of 0.30. The test instead illustrates that an early alert in event-count terms does not necessarily imply low write-byte exposure when individual events transfer large amounts of data.

\subsection{Mixed-Workload Resilience}

A separate mixed-workload experiment evaluated a live BlackBasta binary under four concurrent benign-I/\allowbreak{}O levels: none, light, medium, and heavy. The held-out BlackBasta LOFO checkpoint was used so that BlackBasta remained excluded from model training. One run was conducted at each workload level. The measured alert times were 0.87, 0.55, 0.57, and 0.41 seconds, respectively, and the corresponding cumulative write volumes at the first alert were 302.1, 136.1, 185.8, and 212.8 KiB. No false alerts or ring-buffer event drops were observed during these four runs. These measurements provide preliminary evidence that the held-out checkpoint remained operational under the tested concurrent workloads. Because only one run was performed at each level, the experiment does not establish detection guarantees or production resilience across repeated executions and broader workload conditions. Table 4.12 summarizes the results for the four tested concurrent benign-I/\allowbreak{}O levels.

\begin{table}[htbp]
\centering
\caption{Mixed-workload resilience: BlackBasta (held-out) detection under four concurrent benign I/\allowbreak{}O levels.}\label{tab:4-12}
\begin{adjustbox}{max width=\textwidth}
\small
\begin{tabular}{llllrrr}
\hline
\textbf{I/\allowbreak{}O level} & \textbf{MDR} & \textbf{Detection} & \textbf{WRITE bytes at alert} & \textbf{Files at alert} & \textbf{FP} & \textbf{Dropped} \\
\hline
None & 0.0\% & 0.87 s & 302.1 KiB & 5 & 0 & 0 \\
Light (dd) & 0.0\% & 0.55 s & 136.1 KiB & 6 & 0 & 0 \\
Medium (fio, 4 jobs) & 0.0\% & 0.57 s & 185.8 KiB & 19 & 0 & 0 \\
Heavy (fio + find + cp) & 0.0\% & 0.41 s & 212.8 KiB & 6 & 0 & 0 \\
\hline
\end{tabular}
\end{adjustbox}
\end{table}

Note. Detection is reported as the malware process event-time span at the alert (not wall-clock); the deployed alert configuration (\textmu{}LSTM threshold 0.85, sustained three chunks, then the heuristic filter) was used. Background fio load was time-based to the mounted data volume to preserve the event rate that stresses the ring buffer. Single run per load level on a fresh benign corpus.

\subsection{Limitations}

Three limitations apply to these tests. First, the absence of per-process probability logging means that a zero-alert outcome cannot distinguish a Stage 1 miss from suppression by the sustained-chunk or heuristic layers. Second, T3 and T5 were single-run offline Parquet replay integration tests using a model trained on all fourteen families; they therefore do not support held-out-family or live-execution claims. Third, the benign and mixed-workload experiments used a limited number of workloads and one run per condition. Their results describe operation under the tested conditions but do not estimate production false-positive rates, repeated-run variability, or generalization to broader workload populations.

\section{Aggregate-Feature Baseline Comparison and Threats to Validity}

This section compares the \textmu{}LSTM with aggregate-feature baselines while distinguishing evidence about VFS-derived inputs from claims about classifier architecture. The central contribution of this work is the VFS-level monitoring methodology and the labeled dataset it produces, not the choice of classifier. The stateful \textmu{}LSTM should therefore be read as a practical detector for on-host deployment rather than as a more accurate one.

The aggregate-feature baselines and the \textmu{}LSTM achieved similar mean AUC values on the evaluated VFS dataset: 0.9980 for the Random Forest, 0.9757 for the radial-basis-function SVM, and 0.9752 for the \textmu{}LSTM (Table 4.13). However, their mean MDR values differed: 2.8\%, 2.8\%, and 9.9\%, respectively. The similarity in mean AUC therefore indicates that the twenty-two aggregate VFS features contain substantial discriminative signal, but it does not demonstrate equivalent operational detection performance. In addition, the comparison did not fully control calibration, decision points, stopping rules, or online execution conditions across the three models. The baseline results should consequently be interpreted as an observational comparison rather than evidence that classifier choice has no effect on detection quality.

The similar mean AUC values suggest that classifier sophistication was not the dominant determinant of mean AUC in this experiment. Nevertheless, the lower MDR values of the Random Forest and SVM show that the \textmu{}LSTM did not provide superior detection quality across all metrics. The rationale for retaining the \textmu{}LSTM is therefore its compact streaming state rather than an accuracy or throughput advantage. The recurrent hidden state is fixed in size, but total per-process state also includes a last-read map that grows with the number of distinct paths. A separate matched Stage 1 diagnostic compares this total state and online computation with an incremental Random Forest.

\begin{table}[htbp]
\centering
\caption{Aggregate-feature baselines and the \textmu{}LSTM under the fourteen-fold LOFO evaluation.}\label{tab:4-13}
\begin{adjustbox}{max width=\textwidth}
\footnotesize
\begin{tabular}{lrlllllll}
\hline
\textbf{Family} & \textbf{Malware processes (n)} & \textbf{Write-vol AUC} & \textbf{RF AUC} & \textbf{SVM AUC} & \textbf{\textmu{}LSTM AUC} & \textbf{RF MDR at FPR $\le$ 1\% (\%)} & \textbf{SVM MDR at FPR $\le$ 1\% (\%)} & \textbf{\textmu{}LSTM MDR at FPR $\le$ 1\% (\%)} \\
\hline
Akira & 137 & 0.872 & 1 & 1 & 0.998 & 0 & 0 & 0 \\
Bert & 20 & 1 & 1 & 1 & 1 & 0 & 0 & 0 \\
BlackBasta & 39 & 0.819 & 1 & 1 & 1 & 0 & 0 & 0 \\
Guunra & 29 & 0.120 & 1 & 1 & 1 & 0 & 0 & 0 \\
HelloKitty & 58 & 0.905 & 0.994 & 0.662 & 0.719 & 34.5 & 34.5 & 34.5 \\
IceFire & 35 & 0.997 & 1 & 1 & 1 & 0 & 0 & 0 \\
LockBit & 39 & 1 & 1 & 1 & 1 & 0 & 0 & 0 \\
Monti & 22 & 1 & 1 & 1 & 1 & 0 & 0 & 0 \\
Nitrogen & 35 & 1 & 1 & 1 & 1 & 0 & 0 & 0 \\
RAWorld & 10 & 1 & 1 & 1 & 1 & 0 & 0 & 0 \\
RansomEXX & 38 & 1 & 1 & 1.000 & 1 & 0 & 0 & 0 \\
Royal & 20 & 1 & 1 & 1 & 1 & 0 & 0 & 0 \\
Sodinokibi & 23 & 0.987 & 0.978 & 1.000 & 0.991 & 4.3 & 4.3 & 4.3 \\
SougoLock & 10 & 1 & 1 & 1.000 & 0.945 & 0 & 0 & 100 \\
Mean &  & 0.9072 & 0.9980 & 0.9757 & 0.9752 & 2.77\% & 2.77\% & 9.92\% \\
\hline
\end{tabular}
\end{adjustbox}
\end{table}

Note. The unit in the second column is the number of eligible malware processes, not the number of execution runs. All models use the same held-out-family recording cohorts and the same eligible malware-process sets. The MDR columns are evaluated at a common operating point with a benign FPR of no more than 1\%. The \textmu{}LSTM values in these columns were recalculated from its process-level probabilities and are therefore distinct from the primary Stage 1 MDR values reported at the fixed threshold of 0.30 in Table 4.4. The comparison is observational because the aggregate-feature baselines and the \textmu{}LSTM differ in their state-maintenance, feature-availability, and stopping conditions.

The aligned-cohort comparison produced mean AUC values of 0.9980 for the Random Forest, 0.9757 for the SVM, and 0.9752 for the \textmu{}LSTM. The SVM and \textmu{}LSTM therefore achieved similar mean AUC, whereas the Random Forest achieved a higher mean AUC. At the matched operating point of no more than 1\% benign FPR, the mean per-family MDR was 2.8\% for both aggregate-feature baselines and 9.9\% for the \textmu{}LSTM. The models should consequently not be described as equivalent in overall detection outcomes. The comparison indicates that aggregate VFS features contain substantial discriminative information, but it does not establish that the input representation alone determines performance.

The LOFO comparison in Table 4.13 does not by itself establish deployment equivalence. In a separate matched Stage 1 diagnostic using a common 100-event horizon, the Random Forest and \textmu{}LSTM each detected 18 of 18 replayed malware processes from six families and produced the same MBD3 P75. Median measured per-process detector state was 2,508.2 KiB for the Random Forest and 631.94 KiB for the \textmu{}LSTM. Under the tested heavy-load VPS configuration, the ONNX Random Forest processed 14,902 events/\allowbreak{}s with a 15.2\% event-drop rate, whereas the \textmu{}LSTM processed 9,034 events/\allowbreak{}s with a 58.4\% event-drop rate. Because the replay used all-family deployment models and a small six-family cohort, and the throughput comparison used two repetitions per detector, these findings are supplementary diagnostics rather than LOFO or production-wide estimates. They support a lower measured memory footprint for the \textmu{}LSTM in the evaluated replay cohort, not a universal bounded-memory, accuracy, or throughput advantage.

\section{Chapter Summary}

Chapter 4 presented leave-one-family-out cross-family generalization results on fourteen ransomware families, behavioral-cluster analysis on the reliably detected families, runtime and early-detection measurements from a separately trained prototype deployment model, an error analysis of the residual failure modes across three subsections, ablation studies isolating the contributions of the main design choices, and an end-to-end live-VM validation. The evaluation supports the main claims of the thesis while clearly identifying the conditions under which each claim holds.

The results presented in this chapter provide varying levels of support for the four contributions introduced in Chapter 1. The aggregate-feature comparison in Section 4.9 indicates that the VFS representation carries substantial discriminative signal, while the controlled SougoLock ablation in Section 4.5.3 demonstrates that insufficient representation resolution can cause detection failure within a specific fold. These findings support Contribution C2 as a qualified empirical observation rather than as evidence that representation alone determines detection quality. Contribution C1 is supported by the dataset coverage reported in Table 4.3. Contribution C3 is supported by the behavioral grouping in Section 4.4 and the sensor blind spots identified in Section 4.5.2. Contribution C4 is partially supported by the LOFO results, in which eleven of the fourteen folds achieved MDR below 5\%, together with the model-inference microbenchmark reported in Table 4.10. The remaining failures and deployment limitations prevent these results from being interpreted as universal detection or production-deployment guarantees.

Chapter 5 places these results in a wider context, discusses the specific mitigations suggested by the failure analysis, and outlines the sensor and tokenizer extensions that would address the limitations reported here.

\chapter{Conclusion}

\section{Summary of Findings and Answers to the Research Questions}

This thematic paper investigated the feasibility of ransomware detection on Linux using eBPF-based VFS event streams and a stateful \textmu{}LSTM. The resulting prototype adapts CLEAR-inspired compact tokenization and stateful sequence processing, originally evaluated on Windows NVMe command streams \cite{ref9}, to process-level Linux VFS events comprising read, write, and rename operations.

The primary objective was to evaluate whether the prototype could detect ransomware families excluded from training and to characterize the conditions under which cross-family detection succeeded or failed. The key findings are summarized below.

\subsection{\textmu{}LSTM Classifier Detection}

The LOFO findings provide conditional answers to RQ1, RQ2, and RQ4. The fourteen-fold leave-one-family-out evaluation showed that the selected VFS events and compact tokenization provided useful discriminative information for most evaluated families. Ten folds achieved an MDR of 0.0\%, eleven achieved an MDR below 5\%, and the mean AUC was 0.9752. However, the nine single-binary families limit the interpretation of these results as family-level generalization because their folds also constitute leave-one-binary-out evaluations.

Detection was not uniformly effective across all observed behaviors. SougoLock remained below the fixed Stage 1 threshold because the default token representation did not preserve sufficient offset information, Akira was affected by log-dominant event streams, and HelloKitty produced missed short-lived launcher processes even though its main encryption processes were detected. The fold-specific diagnostic analyses identified plausible limitations and mitigation directions but did not establish universally superior configurations.

\subsection{False-Positive Reduction}

RQ5 is answered as an in-sample analysis, with its trade-offs characterized. The heuristic refinement layer reduced the observed false-positive count from 193 to 4 and increased the mean per-family F1 score from 0.73 to 0.90. However, the benign workloads informed heuristic development, and the diagnostic retained the Stage 1 TP and FN counts rather than re-evaluating true-positive retention. This result should therefore be interpreted as an in-sample false-positive-suppression analysis rather than evidence that false positives can be reduced without increasing missed detections.

\subsection{MBD3 and Data Loss Exposure}

RQ3 concerns the amount of data written before detection. Although MDR indicates whether malware processes were detected, it does not indicate how much data was written before detection. Substantial variation in MBD3 was observed across the held-out-family folds. For example, LockBit and Monti both achieved an MDR of 0.0\%, but their MBD3 values were 103.7 KiB and 16,421.5 KiB, respectively. A low MDR should therefore not be interpreted as uniformly early detection or limited data-loss exposure.

\subsection{Mixed-Workload Feasibility}

The mixed-workload experiment was designed as a supplementary feasibility test rather than a statistical evaluation of workload resilience. BlackBasta served as an encryption-dominant test case for evaluating whether concurrent benign I/\allowbreak{}O prevented event capture or alert generation. The prototype generated an alert at all four evaluated workload levels, with no observed false alerts or ring-buffer event drops. However, because only one run was conducted at each level, the experiment does not quantify run-to-run variability or support generalization to other ransomware families and workload combinations. Repeated trials across multiple families remain necessary for a formal resilience evaluation.

\section{Comparison with the CLEAR Framework}

CLEAR \cite{ref9} reported an MDR of 0.36\% and an MBD3 of 50 MB using Windows NVMe command streams from 47 ransomware families and 137 variants. These values are not directly comparable with the results of This thematic paper. CLEAR defines MDR as missed ransomware traffic volume at an operating point calibrated to one false alarm per 50 GiB of benign traffic, whereas This thematic paper reports the proportion of missed malware processes at a fixed Stage 1 threshold of 0.30.

The two studies also differ in dataset diversity and sensor visibility. CLEAR evaluates a larger collection of families and variants at the NVMe command layer, whereas This thematic paper evaluates fourteen Linux families represented by twenty-eight binaries using process-level VFS events. The prototype also replaces CLEAR's Read-after-Write auxiliary bit with a categorical RENAME flag and uses different offset allocation and overlap-tracking mechanisms.

The results therefore demonstrate the feasibility of adapting CLEAR-inspired compact tokenization and stateful sequence processing to the evaluated Linux VFS environment. They do not establish implementation equivalence with CLEAR or permit direct numerical attribution of performance differences to the operating system, sensor layer, or classifier.

\section{Limitations}

\subsection{Dataset and Environment Scope}

The dataset contains fourteen families and twenty-eight binaries, with nine families represented by a single binary because additional usable samples could not be obtained or detonated successfully. Consequently, these nine folds cannot separate family-level transfer from binary-specific transfer. Although the LOFO test sets are family-disjoint, the training and validation recordings were not grouped by binary hash, victim-file set, or acquisition provenance. All experiments were also conducted on Ubuntu 24.04 LTS in a virtual machine; transfer to other kernels, file systems, Linux distributions, and bare-metal environments was not evaluated.

\subsection{Sensor Visibility and OpCode Coverage}

The sensor monitors selected read, write, and rename paths but does not provide complete visibility into vectored I/\allowbreak{}O, io\_\allowbreak{}uring, memory-mapped writeback, or direct block-device access. Ransomware using uninstrumented paths may therefore evade the current prototype, motivating broader hook coverage in future work.

\subsection{Metric Interpretation}

The reported MDR is process-based and may count short-lived launcher processes as misses even when the associated encryption process is detected, as observed for HelloKitty. MBD3 is calculated only among detected malware processes and should be interpreted as a comparative bytes-to-detection measure rather than an absolute estimate of data loss.

\subsection{Limitations Identified in Live Deployment}

A chunk-zero alert does not necessarily imply negligible damage because substantial write volume may occur within the first one hundred-event chunk. In addition, the reported sub-millisecond latency measures isolated model inference on the Windows development host rather than end-to-end Linux detection and response latency.

\section{Future Work}

\subsection{Expanded Family Coverage}

Future evaluation should expand the number of binaries within each family, include additional Linux distributions, kernels, file systems, and bare-metal hosts, and use hash- or provenance-grouped validation partitions. A disjoint benign workload set is also required to evaluate Stage 2 false-positive suppression and true-positive retention independently.

\subsection{Production Deployment and Sensor Extensions}

Deployment-oriented work should extend sensor visibility to vectored I/\allowbreak{}O, io\_\allowbreak{}uring, and memory-mapped writeback; measure end-to-end latency and resource consumption; and evaluate smaller decision windows for high-throughput ransomware. Persisting per-process Stage 1 and Stage 2 probability traces would further support diagnosis of model misses and heuristic suppression.

\section{Concluding Remarks}

This thematic paper demonstrates that selected Linux VFS events can provide useful behavioral information for process-level ransomware detection using a compact stateful \textmu{}LSTM. Across fourteen held-out-family folds, ten achieved an MDR of 0.0\%, eleven achieved an MDR below 5\%, and the mean AUC was 0.9752. Aggregate-feature baselines achieved similar or higher mean AUC and lower mean MDR under the observational comparison, indicating that the \textmu{}LSTM should be interpreted as a compact streaming design rather than as a universally more accurate classifier.

The results remain conditional on the evaluated families, binaries, sensor hooks, thresholds, and experimental environment. The SougoLock, Akira, and HelloKitty analyses show that detection failures can result from representation resolution, sequence composition, process behavior, and sensor visibility. The study therefore supports the feasibility of CLEAR-inspired stateful sequence processing for Linux VFS monitoring while identifying the additional dataset diversity, sensor coverage, independent benign evaluation, and end-to-end deployment validation required before production effectiveness can be claimed.
